%%%%%%%%%%%%%%%%%%%%%%%%%%%%%%%%%%%%%%%%%%%%%%%%%%%%%%%%%%%%%%%%%%%%%%%%%%%%
% AGUJournalTemplate.tex: this template file is for articles formatted with LaTeX
%
% This file includes commands and instructions
% given in the order necessary to produce a final output that will
% satisfy AGU requirements, including customized APA reference formatting.
%
% You may copy this file and give it your
% article name, and enter your text.
%
% guidelines and troubleshooting are here: 

%% To submit your paper:
\documentclass[draft]{agujournal2019}
\usepackage{url} %this package should fix any errors with URLs in refs.
\usepackage{lineno}
\usepackage[inline]{trackchanges} %for better track changes. finalnew option will compile document with changes incorporated.
\usepackage{soul}
\usepackage{natbib}
\usepackage{siunitx}
\usepackage{booktabs}
\usepackage{multirow}
\usepackage{hhline}
\usepackage[utf8]{inputenc}
\usepackage{textcomp}
\usepackage{amssymb}
\usepackage[T1]{fontenc}
\usepackage{lmodern}
%\UseRawInputEncoding
\setlength\parindent{0pt}
\linenumbers

%%%%%%%
% As of 2018 we recommend use of the TrackChanges package to mark revisions.
% The trackchanges package adds five new LaTeX commands:
%
%  \note[editor]{The note}
%  \annote[editor]{Text to annotate}{The note}
%  \add[editor]{Text to add}
%  \remove[editor]{Text to remove}
%  \change[editor]{Text to remove}{Text to add}
%
% complete documentation is here: http://trackchanges.sourceforge.net/
%%%%%%%

\draftfalse

%% Enter journal name below.
%% Choose from this list of Journals:
%
% JGR: Atmospheres
% JGR: Biogeosciences
% JGR: Earth Surface
% JGR: Oceans
% JGR: Planets
% JGR: Solid Earth
% JGR: Space Physics
% Global Biogeochemical Cycles
% Geophysical Research Letters
% Paleoceanography and Paleoclimatology
% Radio Science
% Reviews of Geophysics
% Tectonics
% Space Weather
% Water Resources Research
% Geochemistry, Geophysics, Geosystems
% Journal of Advances in Modeling Earth Systems (JAMES)
% Earth's Future
% Earth and Space Science
% Geohealth
%
% ie, \journalname{Water Resources Research}

\journalname{Journal of Geophysical Research: Machine Learning and Computation}


\begin{document}

%%%%%%%%%%%%%%%%%%%%%%%%%%%%%%%%%%%%%%%%%%%%%%%
%  TITLE
%
% (A title should be specific, informative, and brief. Use
% abbreviations only if they are defined in the abstract. Titles that
% start with general keywords then specific terms are optimized in
% searches)
%
%%%%%%%%%%%%%%%%%%%%%%%%%%%%%%%%%%%%%%%%%%%%%%%

% Example: \title{This is a test title}

% \title{High-resolution glacier mass balance predictions using a Long Short-Term Memory network, a case study over the Swiss Alps}

\title{Towards a Global Data-Driven Glacier Mass Balance Model Using Transformers}

%%%%%%%%%%%%%%%%%%%%%%%%%%%%%%%%%%%%%%%%%%%%%%%
%
%  AUTHORS AND AFFILIATIONS
%
%%%%%%%%%%%%%%%%%%%%%%%%%%%%%%%%%%%%%%%%%%%%%%%

% Authors are individuals who have significantly contributed to the
% research and preparation of the article. Group authors are allowed, if
% each author in the group is separately identified in an appendix.)

% List authors by first name or initial followed by last name and
% separated by commas. Use \affil{} to number affiliations, and
% \thanks{} for author notes.
% Additional author notes should be indicated with \thanks{} (for
% example, for current addresses).

% Example: \authors{A. B. Author\affil{1}\thanks{Current address, Antartica}, B. C. Author\affil{2,3}, and D. E.
% Author\affil{3,4}\thanks{Also funded by Monsanto.}}

\authors{Marijn van der Meer\affil{1,2}*, Harry Zekollari\affil{3,1}, Alban Gossard\affil{5}, Jordi Bolibar\affil{5}, Kamilla Hauknes Sjursen\affil{6}, Marin Kneib\affil{1,2,4}, Matthias Huss\affil{1,2,4}, Daniel Farinotti\affil{1,2}}


\affiliation{1}{Laboratory of Hydraulics, Hydrology, and Glaciology (VAW), ETH Zurich, Zurich, Switzerland}
\affiliation{2}{Swiss Federal Institute for Snow and Avalanche Research (SLF), Sion, Switzerland}
\affiliation{3}{Department of Water and Climate, Vrije Universiteit Brussel, Brussels, Belgium}
\affiliation{4}{Department of Geosciences, University of Fribourg, Fribourg, Switzerland}
\affiliation{5}{Univ. Grenoble Alpes, CNRS, IRD, G-INP, Institut des Geosciences de l'Environnement, Grenoble, France}
\affiliation{6}{Department of Civil Engineering and Environmental Sciences, Western Norway University of Applied Sciences, Sogndal, Norway}
%(repeat as many times as is necessary)


% Corresponding author mailing address and e-mail address:

% (include name and email addresses of the corresponding author.  More
% than one corresponding author is allowed in this LaTeX file and for
% publication; but only one corresponding author is allowed in our
% editorial system.)

% Example: \correspondingauthor{First and Last Name}{email@address.edu}

\correspondingauthor{Marijn van der Meer}{vmarijn@ethz.ch}



%%%%%%%%%%%%%%%%%%%%%%%%%%%%%%%%%%%%%%%%%%%%%%%
% KEY POINTS
%%%%%%%%%%%%%%%%%%%%%%%%%%%%%%%%%%%%%%%%%%%%%%%
%  List up to three key points (at least one is required)
%  Key Points summarize the main points and conclusions of the article
%  Each must be 140 characters or fewer with no special characters or punctuation and must be complete sentences

% Example:
% \begin{keypoints}
% \item	List up to three key points (at least one is required)
% \item	Key Points summarize the main points and conclusions of the article
% \item	Each must be 140 characters or fewer with no special characters or punctuation and must be complete sentences
% \end{keypoints}

\begin{keypoints}
\item The Transformer-based Mass Balance Machine generalizes better than LSTM when transferring glacier mass balance models to new regions.
\item Pooling training data and fine-tuning with local observations improves predictions for climatically distant regions.
\item  This first step toward a global model leaves scaling to worldwide data and testing beyond the Northern Hemisphere to future work.
\end{keypoints}


%%%%%%%%%%%%%%%%%%%%%%%%%%%%%%%%%%%%%%%%%%%%%%%
%
%  ABSTRACT and PLAIN LANGUAGE SUMMARY
%
% A good Abstract will begin with a short description of the problem
% being addressed, briefly describe the new data or analyses, then
% briefly states the main conclusion(s) and how they are supported and
% uncertainties.

% The Plain Language Summary should be written for a broad audience,
% including journalists and the science-interested public, that will not have 
% a background in your field.
%
% A Plain Language Summary is required in GRL, JGR: Planets, JGR: Biogeosciences,
% JGR: Oceans, G-Cubed, Reviews of Geophysics, and JAMES.
% see http://sharingscience.agu.org/creating-plain-language-summary/)
%
%%%%%%%%%%%%%%%%%%%%%%%%%%%%%%%%%%%%%%%%%%%%%%%

%% \begin{abstract} starts the second page

\begin{abstract}

Glacier mass balance models that generalize beyond well-monitored regions are critical for global assessments and projections of glacier change. Yet no such globally applicable, fully data-driven model currently exists. Building on the Mass Balance Machine (MBM) framework, previously applied within individual regions (Norway and Switzerland), this study moves towards a multi-region model based on a Transformer architecture. Unlike the Long Short-Term Memory (LSTM) network previously used within MBM, which processes months sequentially, the Transformer uses self-attention to simultaneously consider all months of the year and learn their relative importance for point mass balance (PMB) prediction. Driven by monthly climate and static topographic predictors, the Transformer can be applied at any glacier location to reconstruct distributed, monthly mass balance. Using approximately 64,000 PMB measurements from 180 glaciers across six Randolph Glacier Inventory regions, we combine the Transformer with multi-region training and fine-tuning to evaluate its ability to generalize across climatically and topographically diverse settings. For within-region applications, the Transformer matches or outperforms the LSTM baseline and consistently surpasses it in cross-regional transfer. Aggregating training data across four climatically diverse source regions further enhances zero-shot (without fine-tuning) performance for most target regions, particularly the most distant ones. Fine-tuning the multi-region model with only 10\% of local observations reduces RMSE by up to 73\% for the most challenging targets. These findings represent an initial step toward a global, data-driven glacier mass balance model\change{ that integrates a transferable architecture, multi-region training, and efficient adaptation to local topo-climatic glacier conditions.}{, identifying a transferable architecture, multi-region training, and efficient local adaptation as key components.}



\end{abstract}

\section*{Plain Language Summary}
Glaciers worldwide are losing mass at an accelerating rate, but tracking how much ice is gained or lost each year requires direct measurements that exist for only a small fraction of glaciers globally. Most well-monitored glaciers are concentrated in a few regions, such as the Alps and Scandinavia, leaving much of the world poorly observed. We developed a machine-learning model, based on a Transformer (an artificial-intelligence method originally designed for language processing), that estimates glacier mass gain and loss from climate and terrain data alone. Unlike earlier versions trained on a single well-monitored region, our model combines data from four climatically different regions to be more broadly applicable. We then tested how well the trained model predicts mass balance in regions it had never seen, including Alaska, Svalbard, and Central Asia. Predictions made without any local data were often reasonably accurate, and adding just a small amount of local monitoring data, only 10\% of a typical dataset, further reduced errors by up to 73\% in the most challenging regions. These results suggest that even modest new monitoring efforts in poorly observed regions could substantially improve global estimates of glacier change, an important step toward a truly global glacier model.

\section{Introduction}\label{sec:introduction}
% General intro (compressed, global framing):
Glaciers worldwide have experienced accelerated mass loss over recent decades \citep[e.g.,][]{Jakob_2023, Dussaillant_2025, GlaMBIE_2025} and are projected to continue losing mass throughout the 21st century and beyond \citep[e.g.,][]{Rounce_2023, Zekollari_2024, Schuster_2025}. The annual number of glaciers disappearing is expected to peak at up to 4,000 around the mid-21st century, with the timing and magnitude varying significantly by regional climate and glacier size \citep{VanTricht2026}. These losses have substantial environmental and societal consequences \citep{IPCC2023}, highlighting the need for robust and broadly applicable models of glacier change. 

% MB definition + traditional models + PMB scarcity (compressed):
\add{Glacier mass balance comprises surface, internal, basal, and frontal components} \citep{Cogley_2011}\add{. Where tidewater glaciers do not predominate and strong geothermal heat flux is absent, the latter three are small, and the temporal evolution of glacier volume is primarily governed by the surface mass balance, defined as the net sum of surface accumulation and ablation} \citep{Cogley_2011, Cuffey_and_Paterson2010}. \remove{Glacier mass balance, defined as the net sum of accumulation and ablation, is the primary process-level variable governing the temporal evolution of glacier volume.} Its spatial distribution with elevation establishes \add{surface} mass balance as the most direct link between a glacier and its climatic context \citep{Oerlemans_2001, Cogley_2011, Benn_2014}. Physically based models, including temperature-index and energy-balance formulations \citep[e.g.,][]{Hock_2005, Huss_2008, Huss_2015}, are the standard approach for estimating mass balance. These models are typically calibrated against geodetic mass balance \citep{Zekollari_2024}, which, since approximately 2000, has provided near-global, glacier-by-glacier coverage of volume change \citep[e.g.,][]{Dussaillant_2019, Hugonnet_2021, Berthier_2023, GlaMBIE_2025}. However, because geodetic estimates integrate mass change over several years into a single value, multiple physically distinct combinations of parameters can reproduce the same observed volume change. The resulting total may therefore be correct even when its underlying spatial and seasonal distribution is not, a limitation known as \textit{equifinality} \citep{Finger_2011, Gabbi_2014, Rounce_2020, Cremona_2026}. Addressing this limitation requires independent, seasonally resolved measurements, most commonly in situ point surface mass balance (PMB) data collected at individual stake locations \citep{Radic_2011, Farinotti_2015, Vincent_2017}. These measurements are labor-intensive to obtain and, as a result, are sparse and unevenly distributed. The World Glacier Monitoring Service (WGMS) reports records for approximately 500 of the world's more than 200,000 glaciers, with seasonal, point-scale records for fewer than 160, concentrated in a limited number of well-monitored regions \citep{WGMS2023}. This figure, however, understates the true extent of PMB monitoring, as many national and regional programs collect stake measurements that are never reported to this global database.


% ML models + MBM lineage:
Machine learning (ML) provides a data-driven alternative to physically based models by identifying statistical relationships between climate, topography, and mass balance \change{directly from observational data, rather than relying on predefined physical parameterizations}{directly from observational data, without predefined physical parameterizations}. ML approaches have demonstrated increasing potential for mass balance estimation \citep[e.g.,][]{Bolibar_2022, Anilkumar_2023, Guidicelli_2023, vanderMeer_2025, Maslov_2026}. These ML models address the monitoring imbalance and, because they can be trained on point-scale observations rather than a single glacier-integrated value, provide a means to resolve the spatial and seasonal distribution of mass balance. This aspect remains unconstrained by most large-scale modeling approaches, which typically calibrate parameters against a single, glacier-wide geodetic observation. Within this context, the Mass Balance Machine (MBM) framework was introduced to model PMB directly from climate and topographic predictors. The MBM was initially based on gradient-boosted decision trees (XGBoost; \citealp{Sjursen2025})\add{, which performed competitively with physically based models, including temperature-index approaches,} and was later reformulated using a Long Short-Term Memory (LSTM) network to explicitly represent the temporal evolution of accumulation and ablation within a hydrological year \citep{vanderMeer_2026}\add{. This LSTM-based MBM outperformed both XGBoost and a multilayer perceptron (MLP).}


% The gap - global ML models:
Previous MBM applications were developed and evaluated in single, exceptionally well-monitored regions (Norway; \citealp{Sjursen2025} and Switzerland; \citealp{vanderMeer_2026}) and demonstrated strong generalization to glaciers held out from training. This new study investigates whether such a model, once trained, can also generate reliable predictions in entirely different regions not encountered during training. The initial analysis examines whether a model trained on data-rich regions, a single region or several pooled together, can be directly applied to data-sparse regions without local retraining, a so-called ``zero-shot approach''. If this zero-shot approach does not yield sufficient accuracy, the study evaluates whether fine-tuning with a limited amount of data in the target region can efficiently adapt the model. A concurrent effort by \citet{Maslov_2026} explored a related transferability question using a hybrid physics-machine-learning model, in which a shallow neural network learns spatially and temporally varying corrections to a temperature-index backbone. In contrast, the approach we present here learns climate-to-mass-balance relationships directly from observations, without imposing a physically prescribed structure. This methodological choice is intentional: previous MBM implementations have demonstrated that, even when learning exclusively from data, the model identifies physically plausible climatic drivers and temporal dependencies \citep{Sjursen2025, vanderMeer_2026}. These findings suggest that physical consistency does not necessarily require an explicit physical backbone for PMB prediction. As no globally applicable, fully data-driven glacier mass balance model currently exists, evaluating this progression, from direct transfer to local adaptation, constitutes an essential first exploratory step toward such a model. 


% Architectural feature (Transformer):
With the goal of transferability, our study introduces a novel model architecture for data-driven mass balance modeling. \change{Rather than employing the recurrent processing of LSTM networks, where information is passed sequentially from one time step to the next, our study adopts a Transformer architecture based on self-attention, which enables each time step to access information from all others directly.}{LSTM networks pass information sequentially from one time step to the next} \citep{Gers_2000}\add{. Our study instead adopts a Transformer architecture based on self-attention, which enables each time step to access information from all others directly }\citep{Vaswani2017}. Although originally developed for language modeling, Transformers have been successfully applied to sequential Earth science problems, such as weather forecasting and hydrological prediction \citep[e.g.,][]{Bilgin2021, Chattopadhyay2021, Liu2025}. This approach is advantageous in the present context because self-attention operates over the entire hydrological year, allowing the model \change{to dynamically re-weight the importance of different months, rather than compressing the sequence into a single fixed-size state as in LSTM models.}{to dynamically re-weight the importance of different months, with direct access to every time step of the sequence.} Given that the timing of accumulation and ablation varies systematically between climatic regimes \citep{Oerlemans_2001}, this flexibility is likely to benefit cross-regional transfer, where the model must generalize across diverse seasonal patterns.


Guided by this transferable architecture, this study is designed around three main objectives:
\begin{itemize}
    \item \textbf{The value of the Transformer for transfer learning}: how a Transformer architecture, whose self-attention can flexibly re-weight the seasonal signal across differing climatic regimes, compares to the LSTM baseline in zero-shot transfer between regions.
    \item \textbf{The value of multi-region training}: how aggregating training data across several climatically diverse source regions affects zero-shot transfer performance relative to training on any single source region alone.
    \item \textbf{The value of local fine-tuning}: to what extent fine-tuning with a small fraction of local observations from the target region can close the remaining performance gap for the most challenging target regions.
\end{itemize}


% Roadmap:
The remainder of the paper is organized as follows. Section~\ref{sec:methods} details the study regions, Transformer architecture, training and inference configurations, train/test splitting procedure, and fine-tuning strategies. Section~\ref{sec:results} presents results in order of increasing generalization complexity: within-region performance, single-region zero-shot transfer, multi-region zero-shot transfer, and multi-region fine-tuned transfer. Section~\ref{sec:discussion} analyzes the relationship between domain distance and fine-tuning gain and discusses the broader implications of these findings towards building a global, data-driven glacier mass balance model.


\section{Methods}\label{sec:methods}

% Small intro to each section
This study extends the MBM framework \citep{Sjursen2025}, originally based on an XGBoost architecture, and its LSTM-based adaptation for Switzerland \citep{vanderMeer_2026}, retaining the input pipeline, loss formulation, and training protocol. A Transformer-based architecture (hereafter referred to as the Transformer; see Sec.~\ref{subsec:training_framework}) is introduced as a direct replacement for the LSTM baseline. The Transformer is evaluated across two training configurations and two inference modes, resulting in four experimental settings (Sec.~\ref{subsec:training_configurations}). Train/test splits are determined using a Sinkhorn-distance-based procedure to ensure representative partitions across regions with varying levels of PMB data availability (Sec.~\ref{subsec:splitting}). For target regions with limited data or a large geographic separation from source regions, fine-tuning strategies are evaluated to adapt the pre-trained multi-region model using a small subset of observations from the target region (Sec.~\ref{subsec:finetuning}).

\subsection{Predictors and Target Variable}\label{subsec:predictors_target}

The Transformer takes dynamic and static inputs and predicts seasonal PMB as its target variable. \textbf{Dynamic inputs} ($F_d$) are monthly time-varying climatic predictors that form the sequential input to the Transformer encoder; \textbf{static inputs} ($F_s$) are constant topographic descriptors processed by a parallel MLP branch (Sec.~\ref{subsec:training_framework}). Because in-situ PMB observations are only available at seasonal or annual resolution, the model predicts monthly mass balance and aggregates these estimates to the appropriate temporal scale via a custom monthly loss, a standard component of the MBM framework.


\subsubsection{Target Variable: Point Mass Balance}\label{subsec:study_site_and_data}

This study uses 180 glaciers across six RGI regions in the Northern Hemisphere, totaling roughly 64,000 PMB measurements (annual and winter), with sample sizes ranging from 3 glaciers / 460 measurements (Svalbard) to 63 glaciers / 43,201 measurements (Central Europe; Table~\ref{tab:region_overview} and Fig.~\ref{fig:study-regions}).

These PMB measurements are sourced from regional glacier monitoring programs: GLAMOS (Glacier Monitoring Switzerland) for Switzerland \citep{GLAMOS_point}, GLACIOCLIM (les glaciers, un observatoire du climat) for France \citep{GLACIOCLIM_point}, the Norwegian Water Resources and Energy Directorate (NVE) for Norway \citep{Elvehoy2025}, and the Icelandic GlacierWeb Portal for Iceland \citep{Bjornsson1998, Adalgeirsdottir2020}. For the remaining regions (Italy, Austria, Central Asia, and Alaska), no dedicated regional portal consolidating PMB measurements was identified, and the data were therefore obtained from the World Glacier Monitoring Service (\citealp{WGMS2023}), which aggregates contributions from national and regional monitoring programs worldwide.


% Subdivisions
Some RGI regions are geographically extensive and host glaciers spanning a wide range of climatic and topographic settings, motivating finer subdivisions. The Central European Alps (comprising $>$99\% of RGI region 11) are geographically compact and span four countries (Switzerland, France, Italy, Austria; Fig.~\ref{fig:study-regions}e), whose glaciers, despite often being in close proximity, are monitored by distinct national programs and can differ systematically in measurement protocol. This is reflected in the data volume: Switzerland contributes 30,666 measurements across 30 glaciers, versus 8,129 measurements over 9 glaciers for France, 2,479 measurements over 14 glaciers for Italy, and 1,927 measurements over 10 glaciers for Austria (Table~\ref{tab:region_overview}). We therefore subdivide the Central European Alps by country.

Central Asia (RGI region 13) is far more geographically extensive, with glaciers spreading across large distances and pronounced climatic gradients. Because so few glaciers are monitored across this large area \citep{hoelzle2017re}, subdividing it in a climatically meaningful way is not straightforward, and the literature proposes several different classifications \citep[e.g.,][]{kaab2015brief, bolch2019status, Brun_2019, Barandun_2023}. For simplicity, we adopt the existing RGI O2 sub-region (SR) boundaries, which broadly follow similar lines, separating Central Asia into SR 1--4 (Fig.~\ref{fig:study-regions}f). SR~1 and SR~2 are excluded from the analysis because their measurement counts (187 and 69, respectively; Table~\ref{tab:region_overview}) fall below the minimum threshold of 400 PMB observations we consider necessary for meaningful model training and evaluation, leaving SR~3 and SR~4 in the final dataset. \add{Note that Central Asia SR~4 consists of the eastern and western branches of Urumqi Glacier No.~1, which are delineated as two separate glaciers in the RGI.}

For the remaining regions, we retain only the subset of available data most representative of each RGI region: for RGI~08, only Norway is used, as it accounts for the majority of measurements (Fig.~\ref{fig:study-regions}c); for RGI~07, only Svalbard is used, excluding Jan Mayen as it has no PMB measurements (Fig.~\ref{fig:study-regions}b). Alaska and Iceland require no further subdivision and are used as provided (Fig.~\ref{fig:study-regions}a, d).

 
\begin{figure}
  \centering
  \includegraphics[width=15cm]{fig1_study_regions.png}
  \caption []{\small Locations of glaciers with available point mass balance (PMB) measurements are shown as follows: (a) Alaska (Randolph Glacier Inventory - RGI region 01), 12 glaciers, 1952 to 2025; (b) Svalbard (RGI region 07), 3 glaciers, 2006 to 2023; (c) Norway (RGI region 08), 32 glaciers, 1962 to 2021; (d) Iceland (RGI region 06), 55 glaciers, 1986 to 2024; (e) Central European Alps (RGI region 11), 63 glaciers, 1951 to 2025, in this study subdivided/colored by country (CH, FR, IT, AT); (f) Central Asia (RGI region 13), 15 glaciers, 2006 to 2024, subdivided/colored by RGI sub-region (1 to 4). The number of measurements is provided in Table~\ref{tab:region_overview}.}
  \vspace{-3mm}
  \label{fig:study-regions}
\end{figure}

\subsubsection{Dynamic and Static Inputs}\label{subsec:input-variables}

The Transformer uses the same eight dynamic inputs ($F_d = 8$: 2\,m air temperature, total precipitation, forecast albedo, latent and sensible heat fluxes, downward solar radiation, net thermal radiation, and elevation difference) and three static inputs ($F_s = 3$: aspect, slope, sky-view factor) as the LSTM baseline \citep{vanderMeer_2026}, \change{sourced from ERA5-Land and regional digital elevation model (DEM) products from the Open Global Glacier Model, so that performance differences between architectures are attributable to model design rather than to differences in input information.}{so that performance differences between architectures can be attributed to model design alone. Meteorological inputs are taken from ERA5-Land }\citep{ERA5Land}\add{. Topographic information is derived from regional digital elevation models (DEMs). For Switzerland, we follow }\citet{vanderMeer_2026}\add{ and use the annually updated GLAMOS DEMs} \citep{GLAMOS_vol}\add{ together with aspect and slope from the SGI2016 inventory }\citep{SGI2016}.\add{ For all other regions, we use the Copernicus GLO-30 DEM }\citep{COPDEM}\add{, as provided by the OGGM v1.6 preprocessed glacier directories (RGI v6.2 outlines; }\citealp{Maussion_2019})\add{, in which each glacier's DEM is resampled onto a local grid whose resolution increases with glacier size.}

\subsection{Training Framework}\label{subsec:training_framework}

\textbf{Temporal sequence and padding}: Following \cite{vanderMeer_2026}, \add{each PMB observation is paired with a monthly sequence of meteorological inputs covering its valid observation window. This window is defined individually for each observation from its reported start and end dates, for both winter and annual observations, with each date assigned to the nearest month following} \cite{Sjursen2025}\add{. Winter observations thus typically span 6--8 months, ending with the spring survey in March--May, and annual observations 11--12 months, starting as early as August and ending as late as October of the following year. Because the Transformer requires fixed-length sequences, all samples are placed within a common 15-month window (August to October of the following year), and the months outside the valid window are padded. A binary mask identifies the valid months of each observation. It serves both as a key padding mask, excluding padded months from self-attention, and as a temporal mask, ensuring that only valid months contribute to the loss during training and to the seasonal aggregation during inference.}


\textbf{Architecture}: The Transformer adopts a hybrid architecture combining a Transformer encoder for dynamic inputs and an MLP for static inputs (Fig.~\ref{fig:architecture}a). Both branches are processed in parallel before being fused into a per-timestep representation from which monthly mass balance estimates are derived. As PMB observations exist only at seasonal or annual resolution, the predicted monthly mass balance is aggregated to the appropriate scale as described above. Monthly inputs ($F_d = 8$) are linearly projected to a $D$-dimensional space ($D = 64$) and summed with learned positional embeddings encoding month position within the padded hydrological-year sequence; the resulting sequence passes through $N = 4$ Pre-LayerNorm Transformer encoder layers, which improve training stability on the small datasets typical of regional mass balance studies \citep[e.g.,][]{Xiong_2020}.  \add{Within each layer, multi-head self-attention computes, for every month, a learned weighting over all other months of the sequence, so that information from any month can directly inform the representation of any other. This differs from the LSTM baseline, in which information from an earlier month reaches a later month only through successive recurrent states} (Fig.~\ref{fig:app-diff-transformer-lstm}). Static descriptors ($F_s = 3$) are processed in parallel by a shallow MLP (2 hidden layers of width 128, dropout $= 0.1$), yielding a static representation of dimension $S = 128$, which is repeated across the time axis and concatenated with the encoder output to give a fused representation of dimension $D + S = 192$. A linear head (head dropout $= 0.0$) maps this to per-month PMB estimates, summed over the winter and annual masks $m_w$ and $m_a$, as in \citet{vanderMeer_2026}. \add{As the fusion step is a concatenation and the head is linear, static and dynamic information combine additively, and the self-attention is not conditioned on the static predictors. Interactions between topography and climate are therefore represented mainly through the elevation-difference predictor, which is part of the dynamic input sequence. This design follows the LSTM baseline }\citep{vanderMeer_2026}\add{, so that performance differences can be attributed to the temporal encoding.}

\textbf{Optimization and standardization}: Optimization follows the general protocol of \citet{vanderMeer_2026}, with hyperparameters selected via a randomized grid search with 5-fold cross-validation on the training data. We use the AdamW optimizer \citep{Loshchilov2019} with a learning rate of $1\times10^{-3}$ and weight decay of $1\times10^{-4}$. Training is monitored on a validation set, with the learning rate automatically reduced when validation loss plateaus (ReduceLROnPlateau scheduler), and training is stopped early if no improvement is seen after 20 epochs (maximum 150 epochs). A weighted sampler ensures that winter and annual observations contribute equally to each training batch. All predictors and target PMB values are standardized to zero mean and unit variance using training-set statistics before training; model outputs are rescaled back to physical units (m w.e.) during inference. \add{The search spaces, selected hyperparameters, number of trainable parameters, and training times of both architectures are listed in Table }\ref{tab:hyperparameters}.

\begin{figure}
  \centering
  \includegraphics[width=11cm]{fig2_transformer_architecture.pdf}
  \caption []{\small (a) Model architecture: dynamic monthly predictors are encoded using positional embeddings and processed through $N$ Transformer encoder layers with multi-head self-attention. Static predictors are processed by a parallel multilayer perceptron (MLP) and integrated with the encoder output before a linear head generates per-timestep PMB predictions. (b) Zero-shot transfer: a model trained on a source region is directly applied to a target region without utilizing any local target-region data. (c) Fine-tuning transfer: the pre-trained source-region model is adapted using a small fine-tuning set of local target-region observations, followed by evaluation on a held-out target-region test set.}
  \vspace{-3mm}
  \label{fig:architecture}
\end{figure}

\subsection{Training and Inference Experimental Settings}\label{subsec:training_configurations}

The Transformer is evaluated across two training and two inference configurations: 

\begin{itemize}
    \item \textbf{Training configurations}: Training differs in the data used to fit the model. \textbf{Single-region} models are trained on observations from one source region only (unless otherwise specified, on all of the data). \textbf{Multi-region} models are trained jointly on observations from four source regions (CH, NOR, ISL, FR; Sec.~\ref{subsec:study_site_and_data}), with no explicit region identifier or geographic coordinates provided to the model. The idea is that regional identity is instead encoded implicitly through the meteorological forcing: ERA5-Land climate variables vary continuously across space, so glaciers in climatically similar locations will naturally receive similar inputs regardless of which region they belong to. This means the model learns climate-to-mass-balance relationships that are not region-discrete, though it cannot distinguish two stakes with identical forcing but different geographic settings.
    
    \item \textbf{Inference configurations}: differ in whether the model has seen any data from the target region before prediction (Fig.~\ref{fig:architecture}b,c). In the \textbf{zero-shot} transfer (Fig.~\ref{fig:architecture}b), a trained model is applied directly to a target region with no local adaptation (no local PMB observations from the target region are used at any point), giving a clean measure of how well the learned representations generalize across climatic and glaciological regimes. In the \textbf{fine-tuned} transfer (Fig.~\ref{fig:architecture}c), the pre-trained source-region model is instead adapted to the target region using a limited number (10\%) of local PMB observations (the fine-tuning set) before evaluation on a held-out test set (Sec.~\ref{subsec:finetuning}). Fine-tuning begins from pre-trained weights, preserving the physical relationships already encoded during pre-training while allowing the model to specialize in the target regime.
\end{itemize}


Combining these configurations, we select the following four experimental settings:
\begin{itemize}
    \item \textbf{Within-region}: a single-region model trained and evaluated within the same region on held-out glaciers (Sec.~\ref{subsec:within-region}). This differs from a cross-regional transfer (detailed below) in that no domain shift is involved.
    \item \textbf{Single-region zero-shot}: a single-region model applied directly to a target region (Sec.~\ref{subsec:results-zeroshot}).
    \item \textbf{Multi-region zero-shot}: the multi-region model applied directly to a target region (Sec.~\ref{subsec:results-multi-region}).
    \item \textbf{Multi-region fine-tuned}: the multi-region model adapted to a target region using a small local fine-tuning set of the target region (Sec.~\ref{subsec:results-finetuning}).
\end{itemize}

We do not consider the single-region fine-tuned option, as the multi-region model, having been exposed to a broader range of climatic and topographic regimes during pre-training, provides a more general starting point for adaptation to new regions.


\change{\textbf{Choice of glaciers in multi-region set}}{\textbf{Choice of source regions in multi-region set}}: The four source regions for multi-region training (CH, FR, NOR, ISL) are selected on the basis of both data availability and climatic diversity: they are among the regions with the densest and longest PMB records in the dataset (Table~\ref{tab:region_overview}), and together span a range of climatic regimes, from continental Alpine (CH, FR) to maritime Scandinavian (NOR) and subpolar oceanic (ISL), providing a broad foundation from which to learn generalizable climate-to-mass-balance relationships. The remaining regions serve as held-out targets for zero-shot and fine-tuned transfer (Sec.~\ref{subsec:results-zeroshot}--\ref{subsec:results-finetuning}), and are deliberately chosen to span a range of transfer difficulties. Italy and Austria, situated within the same Alpine range as the Swiss and French source regions, are expected to represent relatively easy transfer targets with limited domain shift. Svalbard, while sharing some characteristics with the Icelandic source region, differs in its high-Arctic setting and sparse monitoring record, representing an intermediate difficulty. Alaska and Central Asia, geographically and climatically distant from the source regions, are expected to pose the hardest generalization challenge. This range of target difficulties allows the model's generalization capability to be assessed across a spectrum from near-interpolation to extrapolation. 


\subsection{Train/Test Splitting for Within-Region Training}\label{subsec:splitting}

% Intro to train/test splitting
The experimental settings described above require the model to generalize to glaciers excluded from the training data. In the within-region setting, evaluation is performed on glaciers withheld from the same region used for training. In contrast, in cross-regional settings, evaluation is conducted on the entire target region, as the model has not been exposed to any data from that region except during fine-tuning (Sec.~\ref{subsec:finetuning}).

In the within-region setting, designing train and test splits that accurately reflect the model's generalization task is essential. An inadequately chosen split, such as one where test glaciers are atypically ``easy'' or climatically unrepresentative, may yield an overly optimistic or pessimistic assessment of the model's generalization capability. More generally, each split defines an empirical distribution of climate and topographic features that the model is expected to learn. A representative split ensures that the training distribution aligns with the evaluation distribution, \change{allowing the model to be tested on interpolation within a known domain rather than extrapolation to an unrepresented one.}{so that the model is tested on interpolation within a known domain.} In the within-region setting, this necessitates selecting train and test partitions with similar feature distributions. 

% GL splits:
Glacier-level splits, in which entire glaciers are withheld from training, are used for within-region applications because they provide a more rigorous assessment of spatial generalization to previously unseen glaciers than measurement-level splits, which only exclude individual stake-year observations from glaciers otherwise included in training (following \citet{Sjursen2025, vanderMeer_2026}). However, random assignment of glaciers to training and testing sets can produce unrepresentative partitions, such as a held-out set concentrated at atypically high or low elevations, potentially resulting in misleading evaluations of generalization performance. To address this, the glacier-level split that minimizes the Sinkhorn distance between the feature distributions of the resulting training and testing subsets is selected. This process is evaluated using global feature scalers fitted across all candidate glaciers in the region for normalization. Section~\ref{subsec:within-region} demonstrates that this informed splitting strategy produces more consistent results than random splits.

% Define Sinkhorn
The Sinkhorn distance \citep[e.g.,][]{Cuturi_2013}, computed through the geomloss package \citep{feydy2019interpolating}, is an entropy-regularized approximation of the Wasserstein (optimal transport) distance between two empirical distributions $\mu$ and $\nu$:
\begin{equation}
    \mathcal{W}_\varepsilon(\mu, \nu) = \min_{\pi \in \Pi(\mu,\nu)} \sum_{i,j} c_{ij}\,\pi_{ij} + \varepsilon\,\sum_{i,j} \pi_{ij} \log \pi_{ij}.
\end{equation}
Here, $\Pi(\mu,\nu)$ denotes the set of joint distributions with marginals $\mu$ and $\nu$, $c_{ij}$ represents the cost of transporting mass from point $i$ to point $j$, and $\varepsilon > 0$ is the entropic regularization parameter. Unlike single-feature summary statistics, the Sinkhorn distance compares the full multivariate joint distribution of climate and topographic features. As a result, two glacier subsets can be considered dissimilar even if they match on mean elevation but differ in the joint distribution of elevation and aspect. This distinction is particularly important in regions with few glaciers and sparse, irregular feature distributions.

% Final splits
Glacier-level splits follow an 80/20 train/test ratio in terms of PMB measurements. To identify the best partition, 100 candidate splits are generated by randomly assigning glaciers to train and test sets, and the split that minimizes the Sinkhorn distance between the resulting feature distributions while falling within 5\% of the target measurement fraction is selected. 


% Closing note, relocated from the old intro paragraph
The Sinkhorn distance is also valuable beyond within-region splitting. In cross-regional settings, it provides a means to quantify the difference between source and target domains. As discussed in Sec.~\ref{sec:disc-sinkhorn}, this distance is related to the potential benefit a target region may gain from fine-tuning.

\subsection{Fine-Tuning}\label{subsec:finetuning}

Zero-shot transfer is effective for regions that are topo-climatically similar to the source; however, performance declines in target regions where the pre-trained model must extrapolate beyond its training distribution. Fine-tuning mitigates this limitation by adapting a pre-trained model from the source region to a limited set of PMB observations from the target region, without requiring re-training from scratch (Fig.~\ref{fig:architecture}c). \add{Conceptually, this is analogous to the regional calibration of physically based mass balance models against local observations.} Two fine-tuning strategies are compared:

\begin{itemize}
    \item \textbf{Head-only}: All pre-trained weights remain fixed, and only the output head\add{, a single linear layer mapping the fused representation to monthly PMB,} is updated. \add{This re-weights the features learned during pre-training without modifying them.} This approach substantially reduces the number of trainable parameters and limits the risk of overfitting when only a small number of local samples are available.
    \item \textbf{Full}: All parameters, including the encoder, \change{static MLP, fusion, and head}{static MLP, and head}, are updated and trained jointly at a reduced learning rate ($1 \times 10^{-4}$). \add{As the fusion step consists of a concatenation of the encoder and static MLP outputs, it has no trainable parameters.} This strategy provides greater adaptation capacity compared to the head-only variant\add{, as it also allows the learned temporal and topographic representations to adapt to the target region}.
\end{itemize}

\textbf{Standardization and optimization}: The mean and standard deviation used to normalize predictors and targets remain fixed at the values estimated from the multi-region source regions. This approach ensures that the fine-tuned model operates within the same input space as the pre-trained model and that learned representations are directly comparable. Standardization also ensures that the Sinkhorn distance is comparable across regions. All strategies are trained for up to 50 epochs with early stopping if no improvement is observed after 10 epochs.

\textbf{Fine-tuning data splits}: For all target regions, available PMB observations are divided into a fine-tuning set and a test set using a random measurement-level split. Individual observations, each corresponding to a unique stake-year combination, are shuffled and assigned independently for winter and annual measurements. Ten percent of measurements are used for fine-tuning, while the remaining ninety percent are reserved for testing. \change{Because individual stake-years, rather than entire glaciers, are assigned at random, every glacier}{Because assignment is done at the stake-year level, every glacier} in the region is represented in both sets. To reduce the influence of a single random split, this partitioning is repeated five times with different random seeds, and mean performance is reported across repetitions. The 10\% fine-tuning fraction is intentionally small to reflect the realistic scenario of a region with very limited observations from the target region. This proportion should be interpreted as \change{a conservative test of the minimum data required for effective adaptation, rather than as an optimized choice.}{a conservative test of the minimum data required for effective adaptation; it was not optimized.} \add{To assess the sensitivity of the results to this choice, fine-tuning was repeated with fractions between 5\% and 50\%, using nested subsets of the local observations. All fractions are evaluated on a common test set, comprising the observations not included in the 50\% subset} (Table~\ref{tab:ft_fraction_counts})\add{, so that differences in performance between fractions can be attributed to the amount of fine-tuning data and not to differences in the test set.}


 
\section{Results}\label{sec:results}
Throughout the first two sections, Transformer performance is compared directly against the LSTM baseline of \citet{vanderMeer_2026}. The two architectures differ primarily in how they process the monthly sequence of climate predictors: \change{the LSTM passes information sequentially from one month to the next, compressing the hydrological year into a single evolving hidden state, whereas the Transformer's self-attention mechanism allows every month to attend directly to every other month, in principle preserving and re-weighting seasonal information more flexibly}{the LSTM passes information sequentially from one month to the next through successive recurrent states, whereas the Transformer's self-attention mechanism allows every month to attend directly to every other month, in principle preserving and re-weighting seasonal information more flexibly} (Sec.~\ref{subsec:training_framework}). This distinction is examined first in the within-region setting, where no domain shift is involved, before assessing its impact under cross-regional transfer.


\subsection{Within-Region Modeling}\label{subsec:within-region}

% Intro first
Since MBM has previously been applied only to Norway \citep{Sjursen2025} and Switzerland \citep{vanderMeer_2026}, its capacity to model glaciers in other regions is first evaluated in the within-region setting (Sec.~\ref{subsec:training_configurations}). In this approach, a single-region model is trained and evaluated within the same region on held-out test glaciers following the glacier-level splitting strategy  (Sec.~\ref{subsec:splitting}). Svalbard (3 glaciers) and Central Asia SR 4 (2 glaciers) are excluded from this evaluation due to an insufficient number of glaciers to form a meaningful partition.

% Performance:
The Transformer demonstrates robust within-region performance across \change{all}{most} evaluated regions, with annual root mean square error (RMSE) values ranging from 0.84 m w.e. in France to 1.58 m w.e. in Alaska, and $R^2_a$ (the coefficient of determination on annual measurements) values between 0.49 and 0.85 for all regions except Central Asia SR 3 (Fig.~\ref{fig:within-reg-transformer}). The Transformer performs comparably to the LSTM baseline \citep{vanderMeer_2026} in all regions (for example, Switzerland: $\mathrm{RMSE_a}$ 1.11 versus 1.03 m w.e. for the LSTM and Transformer, respectively; France: 1.17 versus 0.84 m w.e.; Fig.~\ref{fig:app-within-reg-all}), indicating that the Transformer is at least as effective as the established LSTM baseline for within-region modeling across a range of climatic and topographic conditions.

% Exception of CA 3
Central Asia SR 3 represents a notable exception, demonstrating poor predictive performance for both the Transformer ($\mathrm{RMSE_a} = 2.15$~m~w.e., $R^2_a = -1.94$ for annual PMB; Fig.~\ref{fig:within-reg-transformer}h) and the LSTM (Fig.~\ref{fig:app-within-reg-all}p). The scatter plots reveal a distinct two-cluster pattern in the annual predictions: one cluster is predicted with reasonable accuracy, while the second cluster is not well predicted, corresponding to Turgen-Aksuu and Glacier No.~599 (Kjungei Ala-Too; Fig.~\ref{fig:app_CA3_per_gl}a), both located in the easternmost part of the sub-region (Fig.~\ref{fig:app_CA3_per_gl}b). For these two glaciers, the models consistently generate near-zero estimates, despite the observed values being strongly negative. This result suggests that the characteristics of these glaciers may fall outside the range represented by the training glaciers in this split, or that the predictors for these glaciers lack sufficient variability. As discussed in Sec.~\ref{subsec:study_site_and_data}, Central Asia SR 3 encompasses a broad geographic area with only eight widely dispersed glaciers. Consequently, the difficulty in reproducing the mass balance of these glaciers \change{likely arises from an unrepresentative glacier split rather than a fundamental model limitation.}{likely reflects the small number of widely dispersed glaciers available for training, a data limitation that no choice of split can fully compensate for.}

% Random vs Sinkhorn split: 
The Sinkhorn-based glacier split consistently outperforms random glacier assignment across most regions for both winter and annual PMB (Fig.~\ref{fig:app-random-vs-sinkhorn}), with the Sinkhorn split yielding lower RMSE in the majority of cases. Alaska and France represent partial exceptions, as random splits occasionally perform comparably or better, likely due to the high variability associated with a limited number of test glaciers.

% Summary
Collectively, these results demonstrate that the Transformer model generalizes effectively in the within-region setting, matching or surpassing the LSTM baseline across \change{all}{most} regions where a representative glacier split is achievable. With the exception of Alaska and France, the Sinkhorn-based splitting strategy provides consistently more reliable evaluations than random glacier assignment, thereby validating its use throughout this study.

\begin{figure}
  \centering
  \includegraphics[width=15cm]{fig3_within_region_transformer_only.png}
  \caption []{\small Within-region performance of the Transformer, with a separate model trained and evaluated for each region using held-out test glaciers: Switzerland (a), France (b), Italy (c), Austria (d), Norway (e), Iceland (f), Alaska (g), and Central Asia sub-region (SR) 3 (h). Each panel displays modeled versus observed point mass balance (PMB, m w.e.) for annual (pink) and winter (blue) measurements, with the 1:1 line shown in gray. The root-mean-square error (RMSE, m w.e.), the coefficient of determination (R$^2$) and the mean difference (Bias, m w.e.) are given for the annual ($a$) and the winter ($w$) period.}
  \vspace{-3mm}
  \label{fig:within-reg-transformer}
\end{figure}


%% ------------------------------------
\subsection{Single-Region Zero-Shot Transfer}
\label{subsec:results-zeroshot}
%% ------------------------------------

After establishing that the Transformer generalizes well within diverse regions (Sec.~\ref{subsec:within-region}), the next evaluation focuses on cross-regional zero-shot transfer. In this evaluation, the model is trained on a single source region and applied directly to other regions without any local adaptation (Fig.~\ref{fig:architecture}b), which constitutes a more demanding test of generalization than the within-region setting. Specifically, our analysis examines whether the Transformer generalizes more effectively than the LSTM when the source and target regions differ substantially in their topo-climatic regime. To assess this, both the Transformer and LSTM architectures are trained independently on each of the four source regions (CH, NOR, ISL, FR) and evaluated directly on all remaining regions. 

In nearly all source-target pairs, the Transformer model achieves lower zero-shot RMSE than the LSTM, for both annual and winter RMSE, and across all four source regions (Fig.~\ref{fig:zeroshot_comparison}a-h). To demonstrate the general improvement associated with adopting a Transformer approach, four source-target pairs are highlighted. These pairs represent a range of geographic distances and, by extension, topo-climatic differences, from nearby Alpine transfers to more distant cases such as from Norway to Alaska and from Iceland to Italy (Fig.~\ref{fig:zeroshot_comparison}i-p). Overall, these findings indicate that the Transformer consistently surpasses the LSTM in zero-shot transfer, with this advantage evident across diverse source regions and climatic contexts.

\begin{figure}
  \centering
  \includegraphics[width=15cm]{fig4_zeroshot_comparison.pdf}
  \caption []{\small Performance comparison between the Transformer and LSTM in single-region zero-shot transfer. Panels (a-h) display the RMSE difference ($\Delta$ = LSTM minus Transformer) across all source-target pairs for annual (a-d) and winter (e-h) point mass balance (PMB), grouped by source region. Positive values (blue) indicate superior performance by the Transformer, while negative values (gray) indicate the LSTM performs better. Refer to Fig.~\ref{fig:zs_comparison-absolute} for the corresponding absolute RMSE values. Panels (i-p) present scatter plots of modeled versus observed PMB for four selected source-target pairs, shown separately for the LSTM (top row) and Transformer (bottom row), for annual (pink) and winter (blue) PMB. Reported metrics include annual and winter RMSE, $R^2$, and bias.}
  \vspace{-3mm}
  \label{fig:zeroshot_comparison}
\end{figure}



%% ------------------------------------
\subsection{Multi-Region Zero-Shot Transfer}
\label{subsec:results-multi-region}

After demonstrating that the Transformer outperforms the LSTM in single-region zero-shot transfer (Sec.~\ref{subsec:results-zeroshot}), the next step is to evaluate whether directly aggregating training data from all four source regions (CH, NOR, ISL, FR) when training the Transformer model further enhances zero-shot performance. A single-source model is limited by the climatic and topographic characteristics of its training region, whereas multi-region training broadens this distribution and exposes the model to a wider range of topo-climatic regimes.

In most target regions, the multi-region model matches or surpasses the performance of all single-source models (Fig.~\ref{fig:multi_region_comparison}). Scatter plots for all six target regions and five source configurations are shown in Fig.~\ref{fig:app-multi-region-scatter}, providing a more detailed extension of the bar chart summary. The greatest improvements occur for geographically distant targets. For instance, Central Asia SR 4 achieves an annual RMSE of 0.58~m~w.e. with multi-region training, compared to 1.07 to 3.27~m~w.e. for the four single-source models (Fig.~\ref{fig:multi_region_comparison}f). Alaska also demonstrates improved performance with multi-region training (1.44 versus 1.51 to 2.06~m~w.e.; Fig.~\ref{fig:multi_region_comparison}d). For targets located near at least one source region, such as Italy and Austria, multi-region training \change{performs comparably to the best single-source model (CH) rather than providing further improvement.}{performs comparably to, but not better than, the best single-source model (CH).} These findings suggest that multi-region training does not reduce performance when a geographically proximate source already offers a strong baseline (Fig.~\ref{fig:multi_region_comparison}a, b). 

Two regions deviate from the general pattern in annual RMSE. For Central Asia SR 3, single-source models trained on Iceland and France outperform multi-region training (ISL: 1.38, FR: 1.30, versus 2.03~m~w.e. for the multi-region model; Fig.~\ref{fig:multi_region_comparison}e). For Svalbard, the Iceland single-source model achieves a lower annual RMSE than multi-region training (1.20 compared to 1.34~m~w.e.; Fig.~\ref{fig:multi_region_comparison}c). \change{These exceptions may indicate a trade-off between two approaches to effective transfer: multi-region training expands the topo-climatic envelope encountered during training, but a single source region may still outperform if its climate and topography closely match those of the target, effectively serving as a strong topo-climatic analog. For example, Iceland's subpolar oceanic regime is plausibly a closer analog for Svalbard than the multi-region training distribution, which averages over four climatically diverse source regions.}{These exceptions should not be over-interpreted: the difference is small for Svalbard, and both target regions comprise only a few glaciers (3 and 8, respectively), so that their regional RMSE may largely reflect the characteristics of individual glaciers and not those of the region as a whole. One possible explanation is that a single source region may occasionally provide a closer topo-climatic analog for a given target than the combined training distribution of four climatically diverse regions; Iceland's subpolar oceanic regime, for instance, may be a closer analog for Svalbard. Testing this explanation would, however, require more monitored glaciers within each of these target regions, as well as a larger number of target regions, than available in this study.} The winter comparison broadly mirrors this annual pattern, with some differences in the specific target regions where multi-region training underperforms a single-source analog (Fig.~\ref{fig:app-multi-region-barplot-w}).


% Conclusion and path to fine tuning
The results indicate that multi-region training is not invariably the optimal zero-shot strategy; for certain targets, a carefully selected single-source analog can perform equally well or better. However, identifying the most suitable analog in advance necessitates either prior domain knowledge of the target region or a trial-and-error search among candidate source regions, neither of which is feasible for under-monitored regions. Aggregating regions addresses this limitation by broadening the training distribution, enabling the model to remain competitive across most targets without requiring target-specific selection. Although this approach does not always achieve the performance of the single best analog, it is consistently competitive and, as demonstrated in the following section, can be further enhanced through fine-tuning with a limited amount of local data (Sec.~\ref{subsec:results-finetuning}).



\begin{figure}
  \centering
  \includegraphics[width=13cm]{fig5_zs_pooled_vs_ss_barplot_a.png}
  \caption []{\small Zero-shot annual RMSE (m w.e.) is reported for a multi-region model trained jointly on Switzerland, France, Norway, and Iceland (gray), as well as for four single-region models trained on Switzerland (CH, blue), Norway (NOR, light blue), Iceland (ISL, pink), and France (FR, yellow). These models are evaluated on six target regions: Italy (a), Austria (b), Svalbard (c), Alaska (d), Central Asia sub-region 3 (e), and Central Asia sub-region 4 (f). \change{Except for Central Asia sub-region 3, the multi-region model matches or outperforms all single-region models for most targets, with the greatest advantage observed in geographically distant regions such as Central Asia sub-region 4 and Alaska.}{Except for Svalbard and Central Asia sub-region 3, the multi-region model matches or outperforms all single-region models, with the greatest advantage observed in geographically distant regions such as Central Asia sub-region 4 and Alaska.} Corresponding scatter plots for all targets are presented in Fig.~\ref{fig:app-multi-region-scatter}.}
  \vspace{-3mm}
  \label{fig:multi_region_comparison}
\end{figure}


\subsection{Multi-Region Fine-Tuned Transfer}
\label{subsec:results-finetuning}

Although the multi-region zero-shot model demonstrates competitive performance across most regions, certain targets continue to exhibit high RMSE following multi-region training (e.g., Svalbard: 1.34~m~w.e., Central Asia SR 3: 2.03~m~w.e.; Sec.~\ref{subsec:results-multi-region}). This section evaluates whether incorporating a small subset of PMB observations from the target region can further reduce errors during fine-tuning of the multi-region model, comparing head-only and full fine-tuning strategies (Sec.~\ref{subsec:finetuning}). To achieve this, 10\% of observations from the target region are randomly sampled as a fine-tuning set, with this process repeated five times using independent random seeds to mitigate dependence on a single random draw. Unless otherwise stated, RMSE values reported below are the mean across these five repetitions; the corresponding standard deviations are consistently small (Fig.~\ref{fig:ft_gain}), \change{indicating that fine-tuning performance is robust to the particular subset of local observations sampled, rather than dependent on a favorable or unfavorable random draw.}{indicating that fine-tuning performance is robust to the particular subset of local observations sampled.}

Fine-tuning consistently reduces the RMSE compared to zero-shot performance across all six target regions, with full fine-tuning outperforming head-only adaptation in every instance (Fig.~\ref{fig:ft_gain}). The extent of improvement correlates with the initial zero-shot RMSE: for targets where the multi-region model already performed well (Italy, Austria, Central Asia SR 4; zero-shot RMSE $\leq$ 0.94~m~w.e.), fine-tuning yields only modest additional improvement (for example, Central Asia SR 4: 0.58 to 0.50~m~w.e.; Fig.~\ref{fig:ft_gain}f). In contrast, for targets where zero-shot performance was lower (Svalbard, Central Asia SR 3, Alaska; zero-shot RMSE $\geq$ 1.34~m~w.e.), full fine-tuning substantially reduces the error (Central Asia SR 3: 2.03 to 0.54~m~w.e., Svalbard: 1.34 to 0.77~m~w.e., Alaska: 1.44 to 1.08~m~w.e.; Fig.~\ref{fig:ft_gain}e,c,d). Head-only fine-tuning leaves considerable residual error in these regions, indicating that updating only the output head is insufficient when the target region differs markedly from the source regions. Adapting the encoder's internal representations, as done in the full fine-tuning approach, is therefore necessary to address larger domain shifts. 


\add{These results are robust to the choice of the fine-tuning fraction} (Fig.~\ref{fig:app-ft-fraction})\add{. Evaluated on a common test set for all fractions} (Sec.~\ref{subsec:finetuning})\add{, full fine-tuning reduces the error relative to zero-shot transfer in all target regions already with 5\% of the local observations, and its error continues to decrease, with diminishing returns, as the fraction increases to 50\%. For the most distant targets, most of the gain is achieved with small fractions: in Central Asia SR~3, for instance, the annual RMSE decreases from about 2.0~m~w.e. in zero-shot transfer to about 0.6~m~w.e. at 5\% and 0.4~m~w.e. at 50\%. Full fine-tuning outperforms head-only fine-tuning at all fractions. The training and validation losses at 10\%} (Fig.~\ref{fig:app-ft-loss})\add{ show no sign of overfitting for full fine-tuning, including for Svalbard, where the fine-tuning set comprises only 45 observations: the validation loss decreases and then plateaus without a sustained increase, and early stopping restores the checkpoint with the lowest validation loss.}

Together, these results demonstrate that the multi-region Transformer offers a robust and adaptable initialization for fine-tuning. Incorporating as little as 10\% of observations from the target region substantially reduces errors in regions where zero-shot performance was limited, such as in Central Asia SR 3, whereas regions in which the zero-shot method performed well experience only marginal additional improvement. 

\begin{figure}
\centering
\includegraphics[width=15cm]{fig6_pooled_ft_gain_barplot.png}
  \caption []{\small Annual RMSE (m w.e.) is reported for the multi-region zero-shot baseline (gray), head-only fine-tuning (orange), and full fine-tuning (purple), using 10\% of local PMB observations across six target regions: Italy (a), Austria (b), Svalbard (c), Alaska (d), Central Asia sub-region 3 (e), and Central Asia sub-region 4 (f). Results are presented as mean $\pm$ standard deviation across five independent random splits. Fine-tuning consistently reduces RMSE compared to the zero-shot baseline, with the largest improvements observed for the most geographically distant targets from the training regions, specifically Central Asia sub-region 3 and Svalbard. Full fine-tuning outperforms head-only fine-tuning in all cases.}
  \vspace{-3mm}
  \label{fig:ft_gain}
\end{figure}

\subsection{Glacier-Scale Illustration: Distributed Mass Balance Fields}\label{subsec:application}

Previous sections evaluated model performance using stake observations aggregated across entire target regions (Sec.~\ref{subsec:results-zeroshot}--\ref{subsec:results-finetuning}). The Transformer can also generate distributed mass balance maps \change{by extrapolating over all glacier grid cells, rather than being limited to stake locations.}{by extrapolating from stake locations to all glacier grid cells.} This process produces spatially and temporally resolved mass balance fields for each glacier, analogous to the single-region LSTM demonstration for Swiss glaciers in \citet{vanderMeer_2026}. This section presents these outputs for eight glaciers from the study regions (Fig.~\ref{fig:application}) and highlights the effects of fine-tuning in topo-climatically distant regions identified previously (Sec.~\ref{subsec:results-finetuning}). Corresponding mass balance gradients for all glaciers and regions analyzed in this study are provided in the Appendix: Switzerland (Fig. \ref{fig:app_gradients_CH}), France-Italy-Austria (Fig. \ref{fig:app_gradients_CEU}), Norway-Svalbard (Fig. \ref{fig:app_gradients_NOR}), Iceland (Fig. \ref{fig:app_gradients_ISL}), Central Asia SR 3-4 (Fig. \ref{fig:app_gradients_CA}), and Alaska (Fig. \ref{fig:app_gradients_ALA}).


Four glaciers are situated in regions included in the multi-region training pool (Sec.~\ref{subsec:training_configurations}): Argentiere (France; Fig.~\ref{fig:application}a-c), Aletsch (Switzerland; Fig.~\ref{fig:application}d-f), Breidamerkurjoekull (Iceland; Fig.~\ref{fig:application}g-i), and Nigardsbreen (Norway; Fig.~\ref{fig:application}j-l). For these glaciers, the zero-shot and fine-tuned fields are nearly indistinguishable, both spatially and in their mass balance gradients (Fig.~\ref{fig:application}y-B). The Pearson correlation coefficient between the two fields is r~=~1.00, and both approaches closely reproduce the stake-derived mean PMB (RMSE $\approx$~0.93--1.41~m~w.e.). Across the majority of the elevation range, the modeled gradients closely follow the stake observations. However, the modeled mass balance gradient flattens toward the tongue, producing more conservative ablation magnitudes at the lowest elevations, a pattern also reported for the LSTM in Switzerland \citep{vanderMeer_2026}. These findings confirm, at the scale of individual glaciers, the region-level result of Sec.~\ref{subsec:results-multi-region}: fine-tuning provides minimal additional benefit once a glacier's topo-climatic setting is well represented in the training distribution, although it does not fully resolve the recurring tongue-elevation bias.

The remaining four glaciers are located in regions not represented in the multi-region training pool: Nordenskioeldbreen (Svalbard; Fig.~\ref{fig:application}m-o), Lemon Creek (Alaska; Fig.~\ref{fig:application}p-r), and Ts. Tuyuksuyskiy and Golubin (Central Asia; Fig.~\ref{fig:application}s-x). In these cases, fine-tuning results in a more substantial and spatially coherent correction. The most pronounced improvement is observed for Ts. Tuyuksuyskiy, where the zero-shot field is consistently too negative compared to the stake network (RMSE = 2.57 m w.e., mean PMB = $-3.43$ m w.e.; Fig. \ref{fig:application}s). Fine-tuning shifts the entire mass balance field toward less negative values (RMSE = 0.49 m w.e., mean PMB = $-1.77$ m w.e.; Fig. \ref{fig:application}t), aligning the fine-tuned curve more closely with observed stake means across the full elevation range (Fig. \ref{fig:application}E). Golubin exhibits a similar, though less pronounced, correction (RMSE: 2.07 to 0.93 m w.e.; Fig. \ref{fig:application}v-x). Nordenskioeldbreen and Lemon Creek display smaller but still favorable adjustments (Nordenskioeldbreen: 1.02 to 0.49 m w.e.; Lemon Creek: 0.91 to 0.61 m w.e.; Fig. \ref{fig:application}m--n, p--q), reflecting the more modest region-level fine-tuning gains observed for Svalbard and Alaska compared to Central Asia SR 3 (Sec. \ref{subsec:results-finetuning}). For all four glaciers, the zero-shot versus fine-tuned fields (Fig. \ref{fig:application}o, r, u, x) and the PMB gradient comparison (Fig. \ref{fig:application}C-F) \change{indicate that these are not localized adjustments; rather, the fine-tuned field differs from the zero-shot field across the entire glacier extent.}{indicate that the adjustments are not localized: the fine-tuned field differs from the zero-shot field across the entire glacier extent.} 

Collectively, these eight examples demonstrate that the distributed mass balance fields generated by the multi-region Transformer are consistent with the region-level results presented in previous sections. When a glacier's topo-climatic setting is already included in the training pool, fine-tuning results in minimal changes to the distributed field, except for a modest reduction in the recurring tongue-elevation bias. In contrast, if the setting is not represented, fine-tuning coherently reshapes the field across the entire glacier surface. This distinction is significant \change{as the correction generalizes spatially rather than being limited to the fine-tuning locations.}{as the correction extends spatially well beyond the fine-tuning locations.} Consequently, even a small local dataset can improve the full distributed estimate, not just the point predictions used for calibration.


\begin{figure}
\centering
\includegraphics[width=11cm]{fig7_application.png}
  \caption []{\small Comparison of distributed and elevation-dependent mass balance (MB) products derived from the multi-region Transformer under zero-shot (ZS) and fine-tuned (FT) transfer. Panels (a-x) show distributed annual MB maps for Argentiere (a-c; France), Aletsch (d-f; Switzerland), Breidamerkurjoekull (g-i; Iceland), Nigardsbreen (j-l; Norway), Nordenskioeldbreen (m-o; Svalbard), Lemon Creek (p-r; Alaska), 
  Ts.~Tuyuksuyskiy (s-u; Central Asia sub-region 3) and Golubin (v-x; Central Asia sub-region 3). Distributed maps are shown for 2021. Observed in-situ stake measurements are overlaid on the maps as circles. Each row shares the same colorbar. In the scatter plots, each point represents one grid cell, the dashed line indicates the 1:1 relationship, and Pearson's r quantifies the spatial agreement between the two model-based products. Panels (y-F) compare annual mass balance-elevation gradients derived from the ZS (pink) and FT (blue) predictions for the same glaciers as panels (a-x) from 2000--2024. Gradients are computed from distributed model predictions over the full glacier surface and binned in 100 m elevation intervals. Mean observed stake values within each elevation bin are indicated by circles.}
  \vspace{-3mm}
  \label{fig:application}
\end{figure}



\section{Discussion}\label{sec:discussion}
\subsection{Data Volume Distribution}\label{disc:data-asymmetry}

This study uses 180 glaciers across six RGI regions in the Northern Hemisphere. These regions were selected based on the availability of long-term point mass balance (PMB) monitoring records (Sec.~\ref{subsec:study_site_and_data}), resulting in a deliberately heterogeneous dataset that encompasses temperate alpine, continental high-mountain, and maritime polar and subpolar environments. Although geodetic mass balance data now provide near-global, glacier-specific coverage of volume change since approximately 2000 \citep[e.g.,][]{Dussaillant_2019, Hugonnet_2021, Berthier_2023, GlaMBIE_2025}, seasonally resolved, point-scale PMB records remain unevenly distributed across these regions, ranging from 43,201 measurements in Central Europe to only 460 in Svalbard (Table~\ref{tab:region_overview}). This disparity reflects real monitoring inequalities in seasonal and spatially distributed mass balance data\change{, rather than a sampling artifact. This asymmetry is treated as a feature rather than a limitation, as it motivates the multi-region training and fine-tuning approach developed in this study. In this framework, data-rich source regions are utilized to inform predictions in data-sparse target regions.}{and motivates the multi-region training and fine-tuning approach developed in this study, in which data-rich source regions inform predictions in data-sparse target regions.} Consequently, the study \change{addresses the practical challenge of validating a global glacier mass balance model, where}{mirrors, on a small scale, the practical challenge that any future global glacier mass balance model will face, where} a limited number of well-monitored regions must guide predictions for many regions with only annual, geodetically constrained data available.



% The study therefore mirrors the practical challenge faced by any attempt to develop a global glacier mass balance model: a handful of well-monitored regions must inform predictions for the many that are not.


\subsection{Contributions of the Transformer, Multi-Region Training, and Fine-Tuning}

The results of this study demonstrate that combining the three new features: (i) a Transformer architecture, (ii) multi-region training, and (iii) local fine-tuning, yields a progressively more capable model to obtain distributed glacier mass balance predictions across diverse climatic and glaciological settings. We discuss the contribution of each feature in turn.

\textbf{Value of the Transformer}: The Transformer consistently outperforms the LSTM in cross-regional zero-shot transfer (Fig.~\ref{fig:zeroshot_comparison})\change{. This suggests that the Transformer's global self-attention mechanism, which retains an explicit per-timestep representation across the full hydrological-year sequence rather than compressing history into a fixed hidden state, produces more transferable internal representations. The advantage is most pronounced for transfers involving substantial climatic distance}{, with the largest advantage for transfers involving substantial climatic distance} (for example, Iceland to Italy; Fig.~\ref{fig:zeroshot_comparison})\remove{, consistent with the hypothesis that per-timestep attention enables the model to re-weight the seasonal signal flexibly when the timing of accumulation and ablation differs between source and target}. \add{In the multi-region setting, the source regions differ in the timing and length of their accumulation and ablation seasons, from the long accumulation seasons of the maritime regimes of Norway and Iceland to the shorter and more intense melt seasons of the continental Alps, so that the most informative months depend on the climatic context. In the Transformer, the weight assigned to each month is computed from the input sequence itself, and every month has direct access to every other month of the observation window. In the LSTM, information from earlier months must instead be carried forward through successive recurrent states, and} \citet{vanderMeer_2026}\add{ found an effective memory of about three to four months for the LSTM in Switzerland. When the seasonal timing of a target region differs from that of the source regions, the relevant information may lie beyond this horizon, which may explain why the advantage of the Transformer increases with climatic distance. As the learned attention weights were not analyzed directly, this interpretation remains to be confirmed.}

\add{This advantage should, however, be interpreted within the general limits of data-driven models, which can only generalize reliably to conditions that are, at least in combination, represented in their training data. The Transformer can re-weight and recombine seasonal patterns encountered during training, for instance when the onset or length of the melt season differs between source and target regions, but no architecture can be expected to infer the mass balance response of seasonal regimes that are entirely absent from the training data. In summer-accumulation regimes such as those of monsoon-influenced glaciers on the Tibetan Plateau, for example, cold summer temperatures coincide with solid precipitation, a combination that is much less common in the source regions considered here. For such regimes, we do not expect zero-shot transfer to perform well, regardless of the architecture: reliable predictions will require either training data that cover these regimes or local fine-tuning, which provides the model with the missing information directly.}


\textbf{Value of multi-region training}: Aggregating across climatically diverse source regions improves zero-shot performance for most target regions, particularly the most geographically distant ones. Its practical value lies not in always matching the best single-source analog, but in being robust across all targets without requiring prior knowledge of which source region is most appropriate. This is a significant practical advantage when applying the model to under-monitored regions. Here, we considered a multi-region setup consisting of source regions selected for their data richness and climatic diversity, which could be expanded in a more comprehensive foundation model by aggregating additional well-monitored regions, for instance, including the full Central European Alps (France, Italy, Austria) or Scandinavia as additional sources, progressively broadening the training distribution as more data become available through initiatives such as the WGMS.

\textbf{Value of fine-tuning}: Fine-tuning consistently reduces RMSE across all target regions, with the largest gains where the zero-shot model had its lowest performance. Incorporating 10\% of observations from the target region is sufficient to substantially close the remaining gap for the most challenging targets, with diminishing returns for targets already well-served by zero-shot transfer. This asymmetry is encouraging from a practical standpoint: zero-shot transfer is already competitive for regions close to the source distribution, whereas for distant regions a small monitoring investment unlocks a large performance gain (42\% reduction in RMSE for Svalbard and 73\% for CA SR 3; Fig.~\ref{fig:ft_gain}). \add{Why a small number of local observations improves performance so markedly can be understood by analogy with the calibration of physically based models. In temperature-index models, for example, parameters such as degree-day factors and precipitation correction factors are calibrated regionally, as they absorb both biases in the climate forcing and regional differences in the sensitivity of mass balance to climate} \citep[e.g.,][]{Huss_2015, Rounce_2020}\add{. Fine-tuning plausibly plays a similar role for MBM: local observations allow the model to adjust its learned climate-to-mass-balance relationships to the forcing biases and climatic sensitivities of the target region. Two results are consistent with this interpretation. First, for glaciers outside the training regions, fine-tuning shifts the distributed mass balance field coherently across the entire glacier surface }(Fig.~\ref{fig:application})\add{, as would be expected from a correction of systematic offsets, analogous to a precipitation or temperature bias correction. Second, head-only fine-tuning, which only re-weights the features learned during pre-training, recovers considerably less of the remaining error than full fine-tuning in the most distant regions} (Fig.~\ref{fig:ft_gain})\add{, suggesting that the model's sensitivity to the monthly climate inputs must also be adjusted, analogous to recalibrating degree-day factors.}

\subsection{Fine-Tuning Gain and Domain Distance}\label{sec:disc-sinkhorn}
Throughout this study, we use the Sinkhorn distance as a practical tool for splitting glaciers into representative train/test partitions (Sec.~\ref{subsec:splitting}), which proves to be effective at this glacier-level granularity: within-region modeling performance is generally better with Sinkhorn-selected splits than with random glacier assignments (Sec.~\ref{subsec:within-region}; Fig.~\ref{fig:app-random-vs-sinkhorn}), with Central Asia SR 3 as an exception, where both the random and Sinkhorn-minimizing split could not overcome the limited number of widely dispersed glaciers.

A natural follow-up question is whether domain distance can predict not only transfer difficulty but also the extent to which a model stands to gain from fine-tuning. Fig.~\ref{fig:app_domain_shift} shows the annual RMSE gain from full fine-tuning (vs. zero-shot) against the Sinkhorn distance computed over climate predictors (Fig.~\ref{fig:app_domain_shift}a), topographic predictors (Fig.~\ref{fig:app_domain_shift}b), and the geographic coordinates of stake measurement locations (not used as predictors; Fig.~\ref{fig:app_domain_shift}c). For all regions except Central Asia, a broadly positive relationship is visible across all three distance metrics: regions that are topo-climatically further away from the multi-region source distribution tend to benefit more from fine-tuning, consistent with the results of Sec.~\ref{subsec:results-finetuning} and with the intuition that zero-shot representations are less transferable to distant targets, leaving more room for \change{local adaptation}{fine-tuning} to close the gap.

Central Asia is a notable outlier in this analysis. Central Asia SR 4 shows a very limited fine-tuning gain despite its considerable geographic distance, as the multi-region zero-shot model already performed well in this region (annual RMSE 0.58 m w.e.; Sec.~\ref{subsec:results-multi-region}), leaving little room for further improvement. \add{However, as Central Asia SR~4 consists of a single glacier system }(Sec.~\ref{subsec:study_site_and_data})\add{, this good zero-shot performance might not be representative of the performance across the whole sub-region.} \remove{In contrast, }Central Asia SR 3, which is characterized by moderate climate and topographic distance but among the largest geographic distances, exhibits the greatest fine-tuning gain (Fig.~\ref{fig:app_domain_shift}c). \change{This inconsistency in fine-tuning gain across Central Asia may be partially attributed to the quality of the forcing data.}{One possible explanation for why fine-tuning gain does not follow topo-climatic distance in either Central Asian sub-region is the quality of the forcing data across the region as a whole.} Reanalysis products have been shown to perform inadequately in Central Asia's mountainous regions, particularly regarding precipitation and near-surface air temperature, where substantial overestimation errors and challenges in capturing local effects have been reported \citep{Hu_2016, gummadi2025assessing, Barandun_2023}. Although ERA5-Land offers improvements over previous reanalysis products, the core difficulties in representing precipitation and temperature in areas with complex terrain and sparse station networks likely remain \citep{immerzeel2015reconciling, miao2024understanding}. \change{Consequently, the Sinkhorn distance between source and target feature climatic distributions may primarily reflect}{If so, the climate-based Sinkhorn distances for both sub-regions may partly reflect} discrepancies in the reanalysis and DEM products themselves, instead of the actual climatic differences at the glacier surface. \add{With only two Central Asian target regions, comprising eight glaciers and a single glacier system, respectively, the available data do not allow us to distinguish forcing biases from sampling effects. Distinguishing the two would require a direct comparison of ERA5-Land with local station observations in the region.}

When both the model and the climate forcing require regional adjustment, fine-tuning on observations from the target region \change{can}{may} partially compensate, improving performance beyond what the feature-space distance would predict. \change{This interpretation also suggests}{If correct, this interpretation would also imply} that for Central Asia, the fine-tuned model's good performance may partly reflect implicit bias correction rather than pure domain adaptation, and that further improvement would require better-constrained local forcing data, for instance from weather station observations or high-resolution regional reanalyses.


The generally positive relationship between domain distance and fine-tuning gain, along with the Central Asia exception, \change{indicates that geographic distance between source and target measurement locations serves as a useful, albeit imperfect, predictor of expected fine-tuning benefit.}{suggests that geographic distance between source and target measurement locations may serve as a first indicator of expected fine-tuning benefit, although with only six target regions, this relationship remains tentative.} This is especially relevant when climate distance metrics may be influenced by forcing biases in the target region. The development of more robust distance metrics could further elucidate mass balance transferability gaps and inform strategies to address these within the MBM framework.


\subsection{\add{Extending to Global Coverage: Limitations of Data and Architecture}}
\add{Although the six RGI regions considered here span a deliberately diverse range of climates, all lie in the Northern Hemisphere at mid-to-high latitudes, and tropical, monsoon-influenced, and Southern Hemisphere regimes are not represented. Extending the approach to these and other settings faces limitations related to both the input data and the model architecture, which we discuss below.}

\change{\textbf{Bias in ERA5-Land forcings}}{\textbf{Forcing biases and implicit downscaling}}: \add{ERA5-Land exhibits systematic biases in regions with complex terrain and sparse station networks} \citep{Hu_2016, Barandun_2023}\add{. In the current framework, the coarse-resolution forcing is related to point-scale conditions only implicitly, through the elevation-difference predictor, and no explicit bias correction is applied. Unlike physically based models, which typically calibrate a local correction term, such as a precipitation factor, against observed mass balance} \citep[e.g.,][]{Huss_2015}\add{, MBM therefore has no dedicated mechanism to compensate for systematic local biases in the forcing, other than fine-tuning on local observations. Such biases may constrain performance in regions such as Central Asia, irrespective of the architecture. Satellite-derived precipitation products or higher-resolution regional reanalyses could partly mitigate this, as could an explicit, learned correction of the forcing }(Sec.~\ref{subsec:towards-global}).

\change{\textbf{Calibration on PMB}}{\textbf{Calibration on point mass balance observations}}: \change{The current model relies exclusively on in-situ point mass balance observations for training and evaluation. Integrating geodetic mass balance estimates directly into the loss function as a spatially distributed, multi-annual constraint, complementary to the sparse stake network, could enhance physical consistency and robustness, particularly for long-term cumulative mass change. Additionally, remote-sensing-derived snow-cover products, such as transient snowline altitudes, could provide further complementary constraints. These products are both spatially distributed and available at sub-annual resolution, which may help capture the within-season timing of accumulation and ablation that neither sparse PMB stakes nor multi-annual geodetic estimates alone can resolve. The limitation is evident in}{The current model relies exclusively on in-situ PMB observations for training and evaluation, so that its predictions are only as well constrained as the spatial coverage of the stake network. This limitation is evident in} the recurring flattening of the modeled elevation gradient toward the glacier tongue (Fig.~\ref{fig:application}). This pattern has also been reported for the LSTM in Switzerland \citep{vanderMeer_2026}, where similar tail-flattening behavior was attributed to two factors: sparse stake coverage on the extensive, low-elevation tongues of large glaciers, which limits the training signal in this zone, and a general tendency of regression models trained under a mean-squared-error objective to shrink predictions toward the conditional mean, thereby muting the most extreme accumulation and ablation values regardless of architecture. Since this bias recurs across both architectures, \change{it likely reflects shared statistical and data limitations rather than an architecture-specific shortcoming.}{it is likely not architecture-specific, but reflects shared statistical and data limitations.} \add{Complementary, spatially distributed observations could help constrain the model in such poorly sampled areas }(Sec.~\ref{subsec:towards-global}).


% Glacier type and missing predictors
\add{\textbf{Glacier types and missing predictors}: The predictor set also does not represent processes specific to certain glacier types. Several glaciers in the study regions are surge-type or debris-covered, yet the model receives no information on surge state or debris thickness, and can therefore not distinguish their mass balance from that of otherwise similar glaciers. Adding such predictors is limited by data availability: debris thickness observations exist only for a subset of glaciers} \citep{DebDaB_2025}\add{, and surge histories are documented for a limited number of regions. For marine-terminating glaciers, frontal ablation is outside the scope of a surface mass balance model, but would need to be accounted for when comparing modeled surface mass balance with total geodetic mass change; calving fluxes can, however, only be estimated indirectly from ice velocity, thickness, and terminus change} \citep{kochtitzky2022unquantified}\add{, with ice thickness in particular remaining poorly constrained at the global scale} \citep{Farinotti2019, Millan2022}\add{.}


\change{\textbf{Towards a global model}}{\textbf{Seasonal timing and the fixed sequence}}: \change{The six RGI regions considered here span only a subset of global glacier diversity, with all source regions located in the Northern Hemisphere at mid-to-high latitudes. Extending the approach beyond the Northern Hemisphere was outside the scope of this study, as the}{The} fixed hydrological-year sequence structure (Sec.~\ref{subsec:training_framework}) underlying the padding and temporal masking of both \change{architectures}{the LSTM and Transformer architectures} assumes a single accumulation-ablation cycle aligned with the Northern Hemisphere calendar. Southern Hemisphere glaciers, where this cycle is offset by roughly six months, would require a calendar-shifted version of the same structure, whereas tropical or monsoon-dominated glaciers, with quasi-continuous mass turnover, would require a more fundamental revision.

\subsection{\add{Extending to Global Coverage: Potential Research Directions}}\label{subsec:towards-global}

\add{Taken together, multi-region training and fine-tuning provide a consistent and additive benefit: multi-region training broadens the model's climatic coverage without any local data, and fine-tuning then adapts it to the target regime. For regions without local observations, the multi-region model can thus provide zero-shot estimates, which can be improved by fine-tuning on all available local observations wherever these exist; the 10\% fine-tuning fraction used here was intentionally conservative, reserving the remaining 90\% for independent evaluation. Beyond this, however, our results do not provide a recipe for which, or how many, local observations are needed: the fine-tuning set was drawn at random at the measurement level} (Sec.~\ref{subsec:finetuning})\add{, and in regions with few monitored glaciers, the contribution of individual glaciers is difficult to assess robustly. Instead of a fixed workflow, we therefore outline several complementary research directions, which could be pursued individually or in combination, that could move the present multi-region model towards a more reliable global application.}

\add{\textbf{Geodetic constraints}: Glacier-specific geodetic mass balance estimates are available for virtually all glaciers since about 2000} \citep{Hugonnet_2021, GlaMBIE_2025}\add{ and offer the most direct way to constrain the model where PMB observations are scarce or absent.} \citet{vanderMeer_2026}\add{ found that a glaciologically trained LSTM reproduced independent geodetic mass balance well across Swiss glaciers of all sizes, but showed persistent biases for some glaciers, which they attributed partly to the absence of geodetic constraints during training. Integrating geodetic estimates into training is, however, not straightforward. They constrain multi-annual, glacier-wide change, whereas PMB constrains seasonal, point-scale variability, and a sufficiently flexible model may fit the aggregate geodetic signal through physically implausible combinations of seasonal and spatial patterns, an equifinality problem analogous to that of geodetically calibrated physically based models.}

\add{\textbf{Snow-cover constraints}: Remote-sensing-derived snow-cover products, such as transient snowline altitudes, could help reduce this ambiguity, as they are spatially distributed and available at sub-annual resolution, and thus constrain the within-season timing of accumulation and ablation that multi-annual geodetic estimates cannot resolve} \citep{Cremona_2025}\add{. Using them as training targets is, however, also not straightforward. Snow cover does not measure mass balance directly: the snowline marks the boundary between snow and ice, which relates to cumulative mass balance only through additional assumptions, for instance on the distribution of snow depth and density. Optical retrievals are further affected by cloud cover, cast shadows, and debris, snow can be difficult to distinguish from firn, and fresh summer snowfall can temporarily mask the seasonal snowline. Translating such observations into a quantitative constraint on monthly mass balance therefore requires an explicit observation model that links modeled mass balance to the observed snow cover.}


\add{\textbf{Physics-informed hybrid models}: Hybrid models, which combine a differentiable temperature-index backbone with data-driven corrections} \citep[e.g.,][]{Maslov_2026}\add{, offer a complementary way to constrain the joint learning from these different observation types, and could also provide the explicit, learned correction of forcing biases that the current model lacks. Such an approach is currently being explored within the MBM framework} \citep{gossard2026one}\add{. The added physical structure, however, comes at the cost of flexibility: wherever the true relationship between climate and mass balance deviates from the imposed physical formulation, a more constrained model is less able to fit it. How much physical structure to impose, and how to combine it with the temporal representation of the Transformer, therefore remain open questions.}


\add{\textbf{Uncertainty estimates and proxies of transferability}: Deciding which glaciers or stakes would add the most value to a region, or whether observations from one region could inform another, would require proxies of transferability, that is, measures that indicate, before any local data are used, how far a target lies from the training distribution. The Sinkhorn distances used in this study }(Sec.~\ref{sec:disc-sinkhorn})\add{ are a first step in this direction, but remain imperfect: distances computed over climate predictors may partly reflect forcing biases, and their relationship to fine-tuning gain rests on only six target regions. A complementary proxy could come from predictive uncertainty. Methods such as deep ensembles, quantile regression, or conformal prediction can provide uncertainty estimates that should increase for inputs far from the training distribution, and could thereby identify where additional observations are most needed. Such estimates are also essential for users to assess the reliability of individual predictions in regions without independent observations. As the current model provides point predictions only, developing uncertainty estimates is a priority for future work.}

\add{\textbf{Coupling with ice-dynamics models}: MBM is a diagnostic model that reconstructs past mass balance from observed climate forcing. Coupling it with a differentiable ice-dynamics model} \citep[e.g.,][]{Jouvet_2022, Bolibar_2023}\add{, such that modeled surface mass balance drives changes in glacier geometry, which in turn feed back into subsequent mass balance predictions, would enable projections of glacier evolution. Because both components are differentiable, such a coupling would also allow MBM to be constrained by observations of geometric change, such as surface elevation change or glacier area, which are available with near-global coverage from satellite archives. Several challenges remain: ice thickness, required to initialize the ice-dynamics model, is poorly constrained for most glaciers} \citep{Farinotti2019, Millan2022}\add{; changes in glacier geometry become detectable only over years to decades, much longer than the monthly time scale of MBM; and calibrating both components against the same aggregate geometric observations risks allowing deficiencies in one component to be compensated by the other, again an equifinality problem.}



\section{Conclusion}\label{sec:conclusion}
This study extended the Mass Balance Machine (MBM) framework that was previously developed for Norway \citep{Sjursen2025} and Switzerland \citep{vanderMeer_2026} \change{to be applicable and transferable to other, less monitored regions.}{to explore whether, and to what extent, it can be transferred to other, less monitored regions.} To this end, we introduced a Transformer-based architecture and a multi-region training and fine-tuning strategy, which we evaluated across a set of target regions deliberately chosen to span a range of transfer difficulties, from geographically close (Italy, Austria) to climatically and geographically distant (Alaska, Central Asia).

Three main findings emerge. First, the Transformer matches or outperforms the established LSTM baseline in the within-region setting across nearly all regions, and consistently outperforms it in cross-regional zero-shot transfer, supporting the hypothesis that its ability to flexibly re-weight the seasonal signal via self-attention \change{yields more transferable representations than the LSTM's fixed-size hidden state.}{yields more transferable representations than the LSTM's sequential recurrent processing, at least for seasonal patterns that are represented in the training data.} Second, aggregating training data across four climatically diverse source regions further improves zero-shot performance for most, though not all, target regions, and is particularly valuable for geographically distant targets where no single source region serves as a strong analog; its main practical advantage is robustness across targets without requiring prior knowledge of which source region to choose. Third, fine-tuning the multi-region model with only 10\% of observations from the target region substantially closes the remaining performance gap for the most challenging target regions, reducing RMSE by up to 42\% for Svalbard and 73\% for Central Asia SR 3\add{ at monitored glaciers}. At the scale of individual glaciers, this gain manifests not only as an aggregate reduction in error but as a spatially coherent correction of the full distributed mass balance field, \change{demonstrating}{suggesting} that a modest local monitoring investment can meaningfully improve glacier-wide, rather than only point-scale, mass balance estimates. Throughout, the Sinkhorn distance proved a practical tool both for constructing representative train/test partitions and\change{, with some caveats where the input forcing itself carries regional biases (e.g., Central Asia), for anticipating which regions stand to gain most from fine-tuning.}{, more tentatively, for indicating which regions may stand to gain most from fine-tuning.}

\change{While the current experiments are a step towards a global, generalizable model, which would require training on all available global mass balance data and systematic evaluation across held-out regions worldwide, they demonstrate the core ingredients that such a model would need: a transferable architecture, diverse multi-region training, and efficient local adaptation. The experiment design used here, with held-out target regions spanning a range of transfer difficulties from geographically close to distant, provides a template for further developing and evaluating future, more comprehensive multi-region models.}{These results represent a first, exploratory step towards a global, data-driven glacier mass balance model, demonstrating the value of a transferable architecture, diverse multi-region training, and efficient local adaptation. At the same time, they delineate what such a model cannot yet do: it cannot be expected to transfer reliably to seasonal regimes absent from its training data, such as monsoon-influenced or tropical glaciers, and our results do not yet indicate which local observations are most valuable for fine-tuning. Reaching reliable global application will therefore likely require additional observational constraints, such as geodetic and snow-cover data, possibly combined with physically informed model structures, as well as uncertainty estimates that indicate where predictions can be trusted. The experiment design used here, with held-out target regions spanning a range of transfer difficulties from geographically close to distant, provides a template for developing and evaluating such future models.}

\newpage
\appendix

\section{Additional Figures and Tables}
\begin{table}[ht]
\centering
\caption{Overview of stake measurement data by RGI region.}
\label{tab:region_overview}
\begin{tabular}{lllrr}
\toprule
RGI Region & RGI ID & Sub-region & Glaciers & PMB measurements (annual and winter) \\
\midrule
\multirow{4}{*}{Alaska} & \multirow{4}{*}{01}
  & 2 & 8 & 666 \\
  & & 4 & 2 & 502 \\
  & & 6 & 2 & 608 \\
  & & \textbf{Total} & \textbf{12} & \textbf{1,776} \\
\midrule
Iceland         & 06 & ISL & 55 & 6,824 \\
\midrule
Svalbard        & 07 & SJM & 3  & 460 \\
\midrule
Scandinavia     & 08 & NOR & 32 & 6,200 \\
\midrule
\multirow{5}{*}{Central Europe} & \multirow{5}{*}{11}
  & CH & 30 & 30,666 \\
  & & FR & 9  & 8,129 \\
  & & IT & 14 & 2,479 \\
  & & AT & 10 & 1,927 \\
  & & \textbf{Total} & \textbf{63} & \textbf{43,201} \\
\midrule
\multirow{5}{*}{Central Asia} & \multirow{5}{*}{13}
  & 1 & 3 & 187 \\
  & & 2 & 2 & 69 \\
  & & 3 & 8 & 4,329 \\
  & & 4 & 2 & 638 \\
  & & \textbf{Total} & \textbf{15} & \textbf{5,223} \\
\bottomrule
\end{tabular}
\end{table}

\clearpage
\begin{table}[ht]
\centering
\caption{Hyperparameter search spaces and selected configurations of the LSTM and Transformer. Both architectures were tuned with a randomized search with 5-fold cross-validation on the training data, using identical inputs, loss function, and training protocol. Selected values are given in bold.}
\label{tab:hyperparameters}
\small
\begin{tabular}{llll}
\toprule
 & Hyperparameter & Search space & Selected \\
\midrule
\multirow{6}{*}{LSTM}
 & Hidden size & 64, 96, 128, 160 & \textbf{96} \\
 & Layers & 1, 2 & \textbf{2} \\
 & Dropout & 0.1, 0.2, 0.3 & \textbf{0.2} \\
 & Direction & unidirectional (fixed) & \textbf{unidirectional} \\
 & Static MLP layers/width & 0, 1, 2 / 32, 64, 128 & \textbf{1 / 128} \\
 & Static MLP dropout & 0.1, 0.2, 0.3 & \textbf{0.3} \\
\midrule
\multirow{7}{*}{Transformer}
 & Model dimension $D$ & 32, 64, 128, 256 & \textbf{64} \\
 & Attention heads & 2, 4, 8 & \textbf{8} \\
 & Encoder layers $N$ & 1, 2, 3, 4 & \textbf{4} \\
 & Feed-forward width & 64, 128, 256 & \textbf{256} \\
 & Dropout & 0.0, 0.1, 0.2 & \textbf{0.0} \\
 & Static MLP layers / width & 1, 2 / 64, 128 & \textbf{2 / 128} \\
 & Static MLP dropout & 0.0, 0.1, 0.2 & \textbf{0.1} \\
\midrule
\multirow{3}{*}{Both}
 & Head dropout & 0.0, 0.05, 0.1 & \textbf{0.0} (both) \\
 & Learning rate & $5\times10^{-4}$, $10^{-3}$, $2\times10^{-3}$ & \textbf{$5\times10^{-4}$} (LSTM), \textbf{$10^{-3}$} (Transformer) \\
 & Weight decay & $10^{-5}$, $10^{-4}$ & \textbf{$10^{-5}$} (LSTM), \textbf{$10^{-4}$} (Transformer) \\
\midrule
 & Trainable parameters & & 115,937 (LSTM), 219,905 (Transformer) \\
  & Training time per model & & $\sim$12\,s (ALA) to $\sim$25\,min (CH), similar for both \\
\bottomrule
\end{tabular}
\end{table}

\clearpage
\begin{table}[ht]
\centering
\caption{Number of point mass-balance measurements used for fine-tuning in each target region and fine-tuning fraction. Winter and annual measurements are sampled separately at the given fraction, and smaller subsets are nested within larger ones. Counts are identical across the five random repeats. All fractions are evaluated on the same fixed test set: the measurements not included in the 50\% subset. The 10\% fraction (bold) is the one used in the main experiments.}
\label{tab:ft_fraction_counts}
\begin{tabular}{l r r r r r r r r r}
\toprule
 & & \multicolumn{7}{c}{Fine-tuning fraction} & \\
\cmidrule(lr){3-9}
Region & Total & 5\% & \textbf{10\%} & 15\% & 20\% & 25\% & 30\% & 50\% & Test \\
\midrule
Italy             & 2479 & 123 & \textbf{247} & 371 & 495 & 619  & 742  & 1239 & 1240 \\
Austria           & 1927 &  95 & \textbf{191} & 288 & 384 & 481  & 577  &  963 &  964 \\
Svalbard          &  460 &  22 & \textbf{45}  &  68 &  91 & 114  & 137  &  230 &  230 \\
Alaska            & 1776 &  88 & \textbf{177} & 266 & 355 & 443  & 532  &  887 &  889 \\
Central Asia SR 3 & 4329 & 215 & \textbf{432} & 648 & 865 & 1082 & 1298 & 2164 & 2165 \\
Central Asia SR 4 &  638 &  31 & \textbf{63}  &  95 & 127 & 159  & 191  &  318 &  320 \\
\bottomrule
\end{tabular}
\end{table}



\clearpage
\begin{figure}
  \centering
  \includegraphics[width=11cm]{A1_diff_LSTM_Transformer.pdf}
  \caption []{\small Comparison of how the LSTM and Transformer access the monthly predictor sequence spanning the hydrological year (October to September) to predict point surface mass balance (PMB). (a) The LSTM processes months sequentially, carrying information forward through a one-way hidden state, so that month $t$ only has access to information from preceding months. (b) The Transformer processes all months simultaneously via self-attention, allowing every month to attend directly to every other month regardless of order.}
  \vspace{-3mm}
  \label{fig:app-diff-transformer-lstm}
\end{figure}

\clearpage
\begin{figure}
  \centering
  \includegraphics[width=11cm]{A2_within_region_lstm_vs_transf_mixed.png}
  \caption []{\small Comparison of within-region performance between the LSTM (top) and Transformer (bottom) architectures. One model is trained and evaluated per region on held-out test glaciers for eight study regions: Switzerland (a, e), France (b, f), Italy (c, g), Austria (d, h), Norway (i, m), Iceland (j, n), Alaska (k, o), Central Asia sub-region 3 (l, p). Each panel shows modeled versus observed point mass balance (PMB, m w.e.) for annual (pink) and winter (blue) measurements, with the 1:1 line shown in gray.}
  \vspace{-3mm}
  \label{fig:app-within-reg-all}
\end{figure}

\clearpage
\begin{figure}
  \centering
  \includegraphics[width=15cm]{A3_CA3_per_gl.png}
    \caption[]{\small Within-region performance of point mass balance (PMB) modeling by the Transformer for Central Asia sub-region 3, colored by glacier. (a) Modeled versus observed PMB (m w.e.) for annual and winter measurements, with the 1:1 line in gray. (b) Locations of the five test glaciers within Central Asia sub-region 3.}
  \vspace{-3mm}
  \label{fig:app_CA3_per_gl}
\end{figure}

\clearpage
\begin{figure}
  \centering
  \includegraphics[width=15cm]{A4_random_vs_smart_split_transformer.pdf}
    \caption[]{\small RMSE difference ($\Delta\mathrm{RMSE} = \mathrm{RMSE}_\text{random} - \mathrm{RMSE}_\text{Sinkhorn}$) between random and Sinkhorn-informed glacier-level train/test splits for annual (a) and winter (b) PMB, across all within-region study regions. Positive values (blue) indicate that the Sinkhorn split yields lower RMSE than a random split; negative values (gray) indicate the reverse. The Sinkhorn-informed split outperforms random assignment for most regions and both seasons, with Alaska and France the main exceptions. Error bars show the standard deviation across 100 random candidate splits.}
  \vspace{-3mm}
  \label{fig:app-random-vs-sinkhorn}
\end{figure}


\clearpage
\begin{figure}
  \centering
  \includegraphics[width=15cm]{A5_zs_comparison_absolute.png}
  \caption []{\small Absolute zero-shot annual (a-d) and winter (e-h) RMSE (m w.e.) for the LSTM (gray) and Transformer (blue), evaluated on all target regions and grouped by source region: Switzerland (a, e), France (b, f), Norway (c, g), Iceland (d, h). The Transformer achieves lower RMSE than the LSTM for the large majority of source-target pairs; see Fig.~\ref{fig:zeroshot_comparison} for the corresponding $\Delta$RMSE summary.}
  \vspace{-3mm}
  \label{fig:zs_comparison-absolute}
\end{figure}

\clearpage
\begin{figure}
  \centering
  \includegraphics[width=13cm]{A6_zs_pooled_vs_ss_barplot_w.png}
  \caption []{\small Zero-shot winter RMSE (m w.e.) for the multi-region model trained jointly on Switzerland, France, Norway, and Iceland (gray) and four single-region models trained on Switzerland (CH, blue), Norway (NOR, light blue), Iceland (ISL, pink), and France (FR, yellow), evaluated on six target regions. Corresponding scatter plots for all targets are given in Fig.~\ref{fig:app-multi-region-scatter}. }
  \vspace{-3mm}
  \label{fig:app-multi-region-barplot-w}
\end{figure}

\clearpage
\begin{figure}
  \centering
  \includegraphics[width=11cm]{A7_zs_pooled_vs_ss_scatter_all.png}
  \caption[]{\small Modeled versus observed point mass balance (PMB, m w.e.) for zero-shot transfer across all six target regions (rows: Italy, a; Austria, b; Svalbard, c; Alaska, d; Central Asia sub-region 3, e; Central Asia sub-region 4, f) and five source configurations (columns: multi-region, 1; single-region Switzerland, 2; Norway, 3; Iceland, 4; France, 5). Annual (pink) and winter (blue) PMB are shown separately, with the 1:1 line in gray. Metrics shown are annual and winter RMSE, $R^2$, and bias. The multi-region model (column 1) achieves the lowest or near-lowest RMSE for most target regions; see Fig.~\ref{fig:multi_region_comparison} for the corresponding RMSE summary.}
  \vspace{-3mm}
  \label{fig:app-multi-region-scatter}
\end{figure}

\clearpage
\begin{figure}
  \centering
  \includegraphics[width=11cm]{A8_sensitivity_FT.png}
  \caption[]{\small Sensitivity of fine-tuning performance to the fine-tuning fraction. Annual RMSE (m w.e.) on a common test set for head-only (orange) and full (purple) fine-tuning of the multi-region Transformer, as a function of the fraction of local PMB observations used for fine-tuning (5--50\%), for Italy (a), Austria (b), Svalbard (c), Alaska (d), Central Asia sub-region 3 (e), and Central Asia sub-region 4 (f). Points and error bars show the mean and standard deviation across five random splits. The gray line and band show the zero-shot RMSE. The dashed vertical line marks the 10\% fraction used in the main analysis. All fractions are evaluated on a common test set, comprising the observations not included in the 50\% subset (Table~\ref{tab:ft_fraction_counts}), so that differences between fractions reflect the amount of fine-tuning data and not differences in the test set; values at 10\% may therefore differ slightly from those in Fig.~\ref{fig:ft_gain}, which are evaluated on the remaining 90\% of observations.}
  \vspace{-3mm}
  \label{fig:app-ft-fraction}
\end{figure}

\clearpage
\begin{figure}
  \centering
  \includegraphics[width=11cm]{A9_FT_loss_curves.png}
  \caption[]{\small Training (dashed) and validation (solid) loss during head-only (left) and full (right) fine-tuning of the multi-region Transformer with 10\% of the local observations, for each target region. Colors denote the five random splits. Circles mark the checkpoint with the lowest validation loss, which is restored after training. $n_{FT}$ gives the number of fine-tuning observations.}
  \vspace{-3mm}
  \label{fig:app-ft-loss}
\end{figure}





\clearpage
\begin{figure}
  \centering
  \includegraphics[width=17cm]{A10_ds_pooled_gain.png}
  \caption[]{\small Sinkhorn distance versus annual RMSE gain from full fine-tuning over zero-shot for the multi-region model, evaluated on six target regions. Columns show the Sinkhorn distance computed over climate predictors only (a), topographic predictors only (b), and the geographic coordinates (longitude and latitude) of the stake measurement locations (c). Each point corresponds to one target region and is labeled by source $\rightarrow$ target.}
  \vspace{-3mm}
  \label{fig:app_domain_shift}
\end{figure}

\clearpage
\begin{figure}
  \centering
  \includegraphics[width=13cm]{A11_gradients_CH.png}
  \caption[]{\small Comparison of annual mass balance-elevation gradients between zero-shot (ZS, pink) and fine-tuned (FT, blue) transfer for the multi-region Transformer, across all glaciers of Switzerland with stake observations. Gradients are derived from distributed model predictions from 2000-2024 over the full glacier surface and binned in 100 m elevation intervals. Solid lines show the mean gradient and shaded bands show the associated variability for each transfer mode. Mean observed stake values within each elevation bin are indicated by open circles.}
  \vspace{-3mm}
  \label{fig:app_gradients_CH}
\end{figure}

\begin{figure}
  \centering
  \includegraphics[width=13cm]{A12_gradients_CEU.png}
  \caption[]{\small Comparison of annual mass balance-elevation gradients between zero-shot (ZS, pink) and fine-tuned (FT, blue) transfer for the multi-region Transformer, across all glaciers of France, Italy and Austria with stake observations. Gradients are derived from distributed model predictions from 2000-2024 over the full glacier surface and binned in 100 m elevation intervals. Solid lines show the mean gradient and shaded bands show the associated variability for each transfer mode. Mean observed stake values within each elevation bin are indicated by open circles.}
  \vspace{-3mm}
  \label{fig:app_gradients_CEU}
\end{figure}

\clearpage
\begin{figure}
  \centering
  \includegraphics[width=11cm]{A13_gradients_NOR_SJM.png}
  \caption[]{\small Comparison of annual mass balance-elevation gradients between zero-shot (ZS, pink) and fine-tuned (FT, blue) transfer for the multi-region Transformer, across all glaciers of Norway and Svalbard with stake observations. Gradients are derived from distributed model predictions from 2000-2024 over the full glacier surface and binned in 100 m elevation intervals. Solid lines show the mean gradient and shaded bands show the associated variability for each transfer mode. Mean observed stake values within each elevation bin are indicated by open circles.}
  \vspace{-3mm}
  \label{fig:app_gradients_NOR}
\end{figure}

\clearpage
\begin{figure}
  \centering
  \includegraphics[width=13cm]{A14_gradients_ISL.png}
  \caption[]{\small Comparison of annual mass balance-elevation gradients between zero-shot (ZS, pink) and fine-tuned (FT, blue) transfer for the multi-region Transformer, across all glaciers of Iceland with stake observations. Gradients are derived from distributed model predictions from 2000-2024 over the full glacier surface and binned in 100 m elevation intervals. Solid lines show the mean gradient and shaded bands show the associated variability for each transfer mode. Mean observed stake values within each elevation bin are indicated by open circles.}
  \vspace{-3mm}
  \label{fig:app_gradients_ISL}
\end{figure}

\clearpage
\begin{figure}
  \centering
  \includegraphics[width=13cm]{A15_gradients_CA.png}
  \caption[]{\small Comparison of annual mass balance-elevation gradients between zero-shot (ZS, pink) and fine-tuned (FT, blue) transfer for the multi-region Transformer, across all glaciers of Central Asia sub-region 3 and 4 with stake observations. Gradients are derived from distributed model predictions from 2000-2024 over the full glacier surface and binned in 100 m elevation intervals. Solid lines show the mean gradient and shaded bands show the associated variability for each transfer mode. Mean observed stake values within each elevation bin are indicated by open circles.}
  \vspace{-3mm}
  \label{fig:app_gradients_CA}
\end{figure}


\clearpage
\begin{figure}
  \centering
  \includegraphics[width=13cm]{A16_gradients_ALA.png}
  \caption[]{\small Comparison of annual mass balance-elevation gradients between zero-shot (ZS, pink) and fine-tuned (FT, blue) transfer for the multi-region Transformer, across all glaciers of Alaska with stake observations. Gradients are derived from distributed model predictions from 2000-2024 over the full glacier surface and binned in 100 m elevation intervals. Solid lines show the mean gradient and shaded bands show the associated variability for each transfer mode. Mean observed stake values within each elevation bin are indicated by open circles.}
  \vspace{-3mm}
  \label{fig:app_gradients_ALA}
\end{figure}






\clearpage

%%%%%%%%%%%%%%%%%%%%%%%%%%%%%%%%%%%%%%%%%%%%%%%
%
% DATA SECTION and ACKNOWLEDGMENTS
%
%%%%%%%%%%%%%%%%%%%%%%%%%%%%%%%%%%%%%%%%%%%%%%%

\section*{Open Research Section}
Monthly averaged reanalysis data from ERA5-Land \citep{ERA5Land, ERA5} is available at https://doi.org/10.24381/cds.68d2bb30 \citep{ERA5Land_CDS_2022}. The point mass balance dataset for Switzerland used in this study \citep{GLAMOS_point} is available at https://doi.org/10.18750/massbalance.point.2025.r2025. The same holds for the Swiss glacier volume change (\citealp{GLAMOS_vol}; \\ https://doi.org/10.18750/volumechange.2025.r2025) and Swiss glacier-wide mass balance datasets (\citealp{GLAMOS_glacier_wide}; \\https://doi.org/10.18750/massbalance.2025.r2025). Point mass balance data for France are provided by the GLACIOCLIM national observation service \citep{GLACIOCLIM_point} and are available at https://glacioclim.osug.fr/Donnees-des-Alpes. Point mass balance data for Norway are provided by the Norwegian Water Resources and Energy Directorate (NVE; \citealp{Elvehoy2025}) and are available at https://doi.org/10.58059/sjse-6w92. Point mass balance data for Iceland \citep{Bjornsson1998, Adalgeirsdottir2020} were obtained from the Icelandic GlacierWeb Portal (https://icelandicglaciers.is/\#/page/map). Point mass balance data for Italy, Austria, Central Asia, and Alaska were obtained from the Fluctuations of Glaciers database of the World Glacier Monitoring Service (WGMS; \citealp{WGMS2023}), available at https://doi.org/10.5904/wgms-fog-2023-09. This study uses the Mass Balance Machine (MBM; https://github.com/ODINN-SciML/MassBalanceMachine), whose source code and data are available on Zenodo (v.0.1.1; \citealp{MassBalanceMachine_2026}). Glacier outlines and DEMs for all regions except Switzerland (Copernicus GLO-30; \citealp{COPDEM}) were obtained from the OGGM v1.6 preprocessed glacier directories (2025.6 release; https://cluster.klima.uni-bremen.de/\textasciitilde \\oggm/gdirs/oggm\_v1.6/L1-L2\_files/2025.6/elev\_bands\_w\_data/; \citealp{Maussion_2019}).


\section*{Conflict of Interest disclosure}
The authors declare that there are no conflicts of interest for this manuscript.


\acknowledgments
HZ was supported by the European Research Council (ERC) under the European Union's Horizon Framework research and innovation program (Grant agreement No.~101115565; ICE\textsuperscript{3} project) and by the Research Foundation---Flanders (FWO) through an Odysseus Type II project (Grant G0DCA23N; ``GlaciersMD'' project). AG acknowledges funding from the MIAI cluster and Agence Nationale de la Recherche (ANR) in the context of France 2030 (grant ANR-23-IACL-0006). The authors thank the Service National d'Observation GLACIOCLIM (program of CNRS-INSU, OSUG, IRD, INRAE, IPEV, M\'et\'eo France) for their glaciological measurement data.



%%%%%%%%%%%%%%%%%%%%%%%%%%%%%%%%%%%%%%%%%%%%%%%
% REFERENCES and BIBLIOGRAPHY
\clearpage 
\bibliography{bibfile.bib} 

% \begin{thebibliography}{}


% \end{thebibliography}


\end{document}